% Options for packages loaded elsewhere
\PassOptionsToPackage{unicode}{hyperref}
\PassOptionsToPackage{hyphens}{url}
\documentclass[
]{article}
\usepackage{xcolor}
\usepackage[margin=1in]{geometry}
\usepackage{amsmath,amssymb}
\setcounter{secnumdepth}{-\maxdimen} % remove section numbering
\usepackage{iftex}
\ifPDFTeX
  \usepackage[T1]{fontenc}
  \usepackage[utf8]{inputenc}
  \usepackage{textcomp} % provide euro and other symbols
\else % if luatex or xetex
  \usepackage{unicode-math} % this also loads fontspec
  \defaultfontfeatures{Scale=MatchLowercase}
  \defaultfontfeatures[\rmfamily]{Ligatures=TeX,Scale=1}
\fi
\usepackage{lmodern}
\ifPDFTeX\else
  % xetex/luatex font selection
\fi
% Use upquote if available, for straight quotes in verbatim environments
\IfFileExists{upquote.sty}{\usepackage{upquote}}{}
\IfFileExists{microtype.sty}{% use microtype if available
  \usepackage[]{microtype}
  \UseMicrotypeSet[protrusion]{basicmath} % disable protrusion for tt fonts
}{}
\makeatletter
\@ifundefined{KOMAClassName}{% if non-KOMA class
  \IfFileExists{parskip.sty}{%
    \usepackage{parskip}
  }{% else
    \setlength{\parindent}{0pt}
    \setlength{\parskip}{6pt plus 2pt minus 1pt}}
}{% if KOMA class
  \KOMAoptions{parskip=half}}
\makeatother
\usepackage{color}
\usepackage{fancyvrb}
\newcommand{\VerbBar}{|}
\newcommand{\VERB}{\Verb[commandchars=\\\{\}]}
\DefineVerbatimEnvironment{Highlighting}{Verbatim}{commandchars=\\\{\}}
% Add ',fontsize=\small' for more characters per line
\usepackage{framed}
\definecolor{shadecolor}{RGB}{248,248,248}
\newenvironment{Shaded}{\begin{snugshade}}{\end{snugshade}}
\newcommand{\AlertTok}[1]{\textcolor[rgb]{0.94,0.16,0.16}{#1}}
\newcommand{\AnnotationTok}[1]{\textcolor[rgb]{0.56,0.35,0.01}{\textbf{\textit{#1}}}}
\newcommand{\AttributeTok}[1]{\textcolor[rgb]{0.13,0.29,0.53}{#1}}
\newcommand{\BaseNTok}[1]{\textcolor[rgb]{0.00,0.00,0.81}{#1}}
\newcommand{\BuiltInTok}[1]{#1}
\newcommand{\CharTok}[1]{\textcolor[rgb]{0.31,0.60,0.02}{#1}}
\newcommand{\CommentTok}[1]{\textcolor[rgb]{0.56,0.35,0.01}{\textit{#1}}}
\newcommand{\CommentVarTok}[1]{\textcolor[rgb]{0.56,0.35,0.01}{\textbf{\textit{#1}}}}
\newcommand{\ConstantTok}[1]{\textcolor[rgb]{0.56,0.35,0.01}{#1}}
\newcommand{\ControlFlowTok}[1]{\textcolor[rgb]{0.13,0.29,0.53}{\textbf{#1}}}
\newcommand{\DataTypeTok}[1]{\textcolor[rgb]{0.13,0.29,0.53}{#1}}
\newcommand{\DecValTok}[1]{\textcolor[rgb]{0.00,0.00,0.81}{#1}}
\newcommand{\DocumentationTok}[1]{\textcolor[rgb]{0.56,0.35,0.01}{\textbf{\textit{#1}}}}
\newcommand{\ErrorTok}[1]{\textcolor[rgb]{0.64,0.00,0.00}{\textbf{#1}}}
\newcommand{\ExtensionTok}[1]{#1}
\newcommand{\FloatTok}[1]{\textcolor[rgb]{0.00,0.00,0.81}{#1}}
\newcommand{\FunctionTok}[1]{\textcolor[rgb]{0.13,0.29,0.53}{\textbf{#1}}}
\newcommand{\ImportTok}[1]{#1}
\newcommand{\InformationTok}[1]{\textcolor[rgb]{0.56,0.35,0.01}{\textbf{\textit{#1}}}}
\newcommand{\KeywordTok}[1]{\textcolor[rgb]{0.13,0.29,0.53}{\textbf{#1}}}
\newcommand{\NormalTok}[1]{#1}
\newcommand{\OperatorTok}[1]{\textcolor[rgb]{0.81,0.36,0.00}{\textbf{#1}}}
\newcommand{\OtherTok}[1]{\textcolor[rgb]{0.56,0.35,0.01}{#1}}
\newcommand{\PreprocessorTok}[1]{\textcolor[rgb]{0.56,0.35,0.01}{\textit{#1}}}
\newcommand{\RegionMarkerTok}[1]{#1}
\newcommand{\SpecialCharTok}[1]{\textcolor[rgb]{0.81,0.36,0.00}{\textbf{#1}}}
\newcommand{\SpecialStringTok}[1]{\textcolor[rgb]{0.31,0.60,0.02}{#1}}
\newcommand{\StringTok}[1]{\textcolor[rgb]{0.31,0.60,0.02}{#1}}
\newcommand{\VariableTok}[1]{\textcolor[rgb]{0.00,0.00,0.00}{#1}}
\newcommand{\VerbatimStringTok}[1]{\textcolor[rgb]{0.31,0.60,0.02}{#1}}
\newcommand{\WarningTok}[1]{\textcolor[rgb]{0.56,0.35,0.01}{\textbf{\textit{#1}}}}
\usepackage{longtable,booktabs,array}
\newcounter{none} % for unnumbered tables
\usepackage{calc} % for calculating minipage widths
% Correct order of tables after \paragraph or \subparagraph
\usepackage{etoolbox}
\makeatletter
\patchcmd\longtable{\par}{\if@noskipsec\mbox{}\fi\par}{}{}
\makeatother
% Allow footnotes in longtable head/foot
\IfFileExists{footnotehyper.sty}{\usepackage{footnotehyper}}{\usepackage{footnote}}
\makesavenoteenv{longtable}
\usepackage{graphicx}
\makeatletter
\newsavebox\pandoc@box
\newcommand*\pandocbounded[1]{% scales image to fit in text height/width
  \sbox\pandoc@box{#1}%
  \Gscale@div\@tempa{\textheight}{\dimexpr\ht\pandoc@box+\dp\pandoc@box\relax}%
  \Gscale@div\@tempb{\linewidth}{\wd\pandoc@box}%
  \ifdim\@tempb\p@<\@tempa\p@\let\@tempa\@tempb\fi% select the smaller of both
  \ifdim\@tempa\p@<\p@\scalebox{\@tempa}{\usebox\pandoc@box}%
  \else\usebox{\pandoc@box}%
  \fi%
}
% Set default figure placement to htbp
\def\fps@figure{htbp}
\makeatother
\setlength{\emergencystretch}{3em} % prevent overfull lines
\providecommand{\tightlist}{%
  \setlength{\itemsep}{0pt}\setlength{\parskip}{0pt}}
\usepackage[]{natbib}
\bibliographystyle{plainnat}
\usepackage{amsmath}
\usepackage{booktabs}
\newcommand{\pkg}[1]{\textbf{#1}}
\newcommand{\proglang}[1]{\textsf{#1}}
\newcommand{\code}[1]{\texttt{#1}}
\newcommand{\K}{\{K\}}
\usepackage{bookmark}
\IfFileExists{xurl.sty}{\usepackage{xurl}}{} % add URL line breaks if available
\urlstyle{same}
\hypersetup{
  pdftitle={: Network-Embedded Strategic Search Simulation in },
  hidelinks,
  pdfcreator={LaTeX via pandoc}}

\title{\textbf{searchnet}: Network-Embedded Strategic Search Simulation
in \textsf{R}}
\author{true}
\date{}

\begin{document}
\maketitle
\begin{abstract}
The NK fitness landscape model is among the most widely used
computational frameworks in strategy and organization theory, yet it
treats the complexity parameter \(K\) as exogenous, the landscape as
fixed, and agents as isolated searchers. The \proglang{R} package
\pkg{searchnet} addresses all three limitations by reformulating the NK
model as a bipartite network of \(M\) actors affiliating with \(N\)
components and embedding it within a Stochastic Actor-Oriented Model
(SAOM). This reformulation reveals that the scalar \(K\) is a special
case of a four-dimensional coupled degree system---actor scope
(\(K_{AC}\)), component popularity (\(K_{CA}\)), actor sociality
(\(K_{AA}\)), and component epistasis (\(K_{CC}\))---whose dynamics are
governed by actors' forward-looking affiliation decisions rather than
imposed by the modeler. The package provides an \proglang{R6}-based
simulation engine built on \pkg{RSiena}, a clean functional API
(\code{saomnk\_env}, \code{saomnk\_model}, \code{saomnk\_run}), coupled
degree visualization through the \(\{K\}\)-4 panel, utility
decomposition diagnostics, exogenous shock infrastructure, and
validation against the classical NK model. Three worked illustrations
demonstrate reproduction of classical NK search dynamics, endogenous
landscape formation with heterogeneous strategies, and
quasi-experimental shock analysis. Benchmarks against \pkg{nkmodel} and
raw \pkg{RSiena} quantify computational overhead and scalability.
\pkg{searchnet} targets researchers in strategic management,
organizational theory, innovation economics, and computational social
science who study search, adaptation, and competitive dynamics on
complex landscapes.
\end{abstract}

\textbf{Keywords}: NK fitness landscape, stochastic actor-oriented
model, bipartite network, strategic search, coupled degree, R.

\section{Introduction}\label{introduction}

The NK fitness landscape model
\citep{Kauffman1987, KauffmanWeinberger1989} has become one of the most
productive computational frameworks in strategic management and
organization theory. Since \citet{Levinthal1997} demonstrated that the
NK model could illuminate the tension between exploration and
exploitation on rugged performance landscapes, hundreds of studies have
deployed it to investigate organizational adaptation, modular design,
imitation, and competitive search
\citep{Baumann2019, Rivkin2000, SiggelkowRivkin2005}.

The model's appeal lies in its parsimony. Two parameters---\(N\), the
number of design choices or components, and \(K\), the degree of
epistatic interdependence among them---generate landscapes ranging from
smooth single-peaked surfaces (\(K = 0\)) to highly rugged terrains
(\(K = N - 1\)) populated by many local optima. This parsimony has also
become the model's principal liability. Four limitations recur across
the literature.

First, the fitness landscape is \textbf{exogenous}: it is drawn once
from a random distribution and does not respond to the positions that
agents occupy. Firms competing on the same landscape experience
identical payoff surfaces regardless of how many competitors share each
component or how concentrated competitive attention has become. Second,
\textbf{\(K\) is a scalar} set by the modeler. In empirical
organizational settings, interdependence is neither uniform nor fixed;
some components are tightly coupled while others operate nearly
independently, and coupling structures change as actors reconfigure
their systems. Third, \textbf{multi-agent extensions} typically treat
agents as independent searchers navigating the same landscape, ignoring
the social and competitive structure through which firms observe,
imitate, and crowd one another. Fourth, the NK model provides \textbf{no
statistical framework} for parameter estimation from empirical data,
making it difficult to discipline simulations with observed
organizational behavior.

This paper introduces \textbf{searchnet}, an \textsf{R} package that
addresses all four limitations simultaneously. The package reformulates
the NK model as a bipartite network of \(M\) actors (firms, agents,
decision-makers) affiliating with \(N\) components (technologies,
resources, markets, design choices). By embedding this bipartite
structure within a Stochastic Actor-Oriented Model
\citep[SAOM;][]{Snijders1996, SnijdersVdBuntSteglich2010}, the search
environment becomes a network data generation process whose
topology---and therefore whose fitness landscape---co-evolves with
actors' affiliation decisions.

This reformulation yields three contributions.

\emph{Endogenization.} The fitness landscape is no longer exogenous.
Landscape ruggedness emerges from the structure of the bipartite
network, which is itself shaped by actors' strategic choices. When
multiple actors affiliate with the same components, they create
epistatic coupling through overlapping activity systems; when they
withdraw, coupling dissolves.

\emph{Decomposition.} The scalar \(K\) is revealed as a special case of
a four-dimensional coupled degree system, denoted \(\{K\}\). The four
dimensions---actor scope (\(K_{AC}\)), component popularity
(\(K_{CA}\)), actor sociality (\(K_{AA}\)), and component epistasis
(\(K_{CC}\))---are individually standard network statistics
\citep{Latapy2008, BorgattiEverett1997}, but the analytical leverage of
the \(\{K\}\) framework comes from treating them as a coupled system.
Because all four derive from the same bipartite incidence matrix, a
change in any one dimension mechanically propagates to the others,
exposing the transmission mechanism through which micro-level
affiliation decisions produce macro-level landscape topology.

\emph{Statistical infrastructure.} Because the model is formalized as a
SAOM, it inherits \textbf{RSiena}'s \citep{Ripley2022} well-developed
statistical machinery, including maximum likelihood estimation,
goodness-of-fit testing, and the capacity to be estimated on empirical
network panel data.

\noindent A full-scale empirical application using eight data layers
from the U.S.~airline industry is reported in
\citet{downing2026architecture}.

\begin{figure}
\centering
\includegraphics[width=0.85\linewidth,height=\textheight,keepaspectratio,alt={The \textbackslash\{K\textbackslash\} Interdependence System: four coupled degree dimensions emerging from the bipartite actor-component network. Adapted from Downing (2026a).}]{searchnet-jss-NEW2_files/figures/cd2026_conceptual_architecture.png}
\caption{The \(\{K\}\) Interdependence System: four coupled degree
dimensions emerging from the bipartite actor-component network. Adapted
from Downing (2026a).}
\end{figure}

The package provides an \textsf{R6}-based simulation engine, a
functional user API that abstracts away \textbf{RSiena}'s internal
naming conventions, comprehensive visualization through the \(\{K\}\)-4
panel and 3D phase space, utility decomposition diagnostics, exogenous
shock infrastructure for quasi-experimental designs, multi-wave Monte
Carlo facilities, specification robustness tools, an interactive
strategy game mode, and a classroom teaching module. Eight vignettes
accompany the package: a quick-start introduction, theoretical
foundations, a full simulation walkthrough with strategies and shocks, a
computational experiment design guide, a causal inference tutorial, an
NK validation guide, a Blume--Gibbs convergence tutorial, and a
Professional Development Workshop curriculum.

The remainder of this paper is organized as follows.
Section\textasciitilde2 reviews the NK and SAOM literatures and
positions the contribution. Section\textasciitilde3 presents the SaoMNK
framework: the bipartite formulation, the \(\{K\}\) interpretive
framework, the SAOM objective function, and the rationale for choosing
SAOMs over ERGMs. Section\textasciitilde4 describes the
\textbf{searchnet} package architecture, user API, visualization tools,
diagnostics, animation pipeline, interactive features, and proof
registry. Section\textasciitilde5 provides three worked illustrations
and introduces basin geometry analysis for the depth-width tradeoff.
Section\textasciitilde6 presents computational benchmarks and
scalability analysis. Section\textasciitilde7 concludes.

\section{Background: NK models and SAOMs}\label{sec:background}

\subsection{The NK fitness landscape}\label{the-nk-fitness-landscape}

The NK model \citep{Kauffman1987, KauffmanWeinberger1989} characterizes
performance landscapes with two parameters. A system of \(N\) binary
components (loci, design choices, technologies) generates \(2^N\)
possible configurations. Each component \(d\) makes a fitness
contribution \(f_d\) that depends not only on its own state but also on
the states of \(K\) other components with which it is epistatically
linked. Total fitness is the average contribution across all \(N\)
components. When \(K = 0\), each component contributes independently and
the landscape has a single global peak. As \(K\) increases toward
\(N - 1\), the landscape becomes increasingly rugged, populated by
exponentially many local optima that trap myopic searchers.

\citet{Levinthal1997} introduced the NK model to management as a
formalization of the exploration--exploitation trade-off: agents that
search locally (single-bit mutations) on rugged landscapes quickly find
nearby peaks but rarely discover the global optimum, while agents that
search more broadly (multi-bit jumps) sacrifice short-run performance
for a chance at higher peaks. This insight has been extended to
organizational design \citep{SiggelkowRivkin2005}, imitation
\citep{Rivkin2000}, modular architectures \citep{EthirajLevinthal2004},
and competitive interaction \citep{Lenox2006}.

Despite its generativity, the classical NK model imposes structural
constraints that limit its applicability to settings where search
environments are dynamic, socially embedded, or empirically observable.
The landscape is drawn once and held fixed; \(K\) is treated as a
uniform scalar; agents search independently; and there is no estimation
framework linking the model to data. Several extensions have relaxed
individual constraints---\citet{Gavetti2005} introduced cognitive
search, \citet{SiggelkowRivkin2005} embedded search in organizational
hierarchies, and \citet{Baumann2019} provided a comprehensive
review---but no existing framework addresses all four limitations
simultaneously within a unified statistical model.

\subsection{Stochastic Actor-Oriented
Models}\label{stochastic-actor-oriented-models}

Stochastic Actor-Oriented Models \citep{Snijders1996,
SnijdersVdBuntSteglich2010} provide a continuous-time Markov model for
the co-evolution of social networks and behavior. At each micro-step, a
single actor is selected according to a rate function, evaluates the
marginal utility of each possible tie change according to an objective
function, and chooses probabilistically via a multinomial logit rule.
The objective function takes the general form

\begin{equation}\label{eq:saom-objective}
  f_i(\mathbf{B}) = \sum_k \theta_k \, s_{ik}(\mathbf{B}),
\end{equation}

\noindent where \(s_{ik}(\mathbf{B})\) are network statistics computed
from the current network \(\mathbf{B}\) and \(\theta_k\) are parameters
governing each actor's preferences. The parameters \(\theta_k\) can be
estimated from observed panel data using the method of moments or
maximum likelihood, implemented in the \textbf{RSiena} package
\citep{Ripley2022}.

SAOMs are well established in sociology and political science for
modeling friendship networks, advice networks, and inter-organizational
partnerships. They have been extended to bipartite (two-mode) networks
\citep{Koskinen2012}, where the model governs actors' affiliation
decisions with respect to events, organizations, or other categorical
entities. \textbf{searchnet} exploits this bipartite extension to model
actors affiliating with components in a search environment.

\subsection{\texorpdfstring{Positioning
\textbf{searchnet}}{Positioning }}\label{positioning}

No existing \textsf{R} package provides the capability that
\textbf{searchnet} offers. The \textbf{RSiena} package
\citep{Ripley2022} implements SAOMs but does not include bipartite
search environment simulation with NK-style fitness landscape analysis.
The \textbf{nkmodel} package \citep{Baumann2019nk} implements classical
NK models but not the network-embedded generalization.
\textbf{searchnet} bridges these two traditions, providing a simulation
engine in which the fitness landscape is an emergent property of actors'
network positions rather than an exogenously imposed parameter.

\section{The SaoMNK framework}\label{sec:framework}

\noindent \textbf{Notation.} Throughout this paper we adopt the notation
of \citet{downing2026architecture}, where the formal SAOM-NK proofs are
developed in full. The bipartite incidence matrix is
\(\mathbf{B} \in \{0,1\}^{M \times N}\) with entries \(b_{ij}\); the NK
configuration vector for a single agent is \(\mathbf{x} \in \{0,1\}^N\);
the influence matrix is \(\mathbf{E}\) (implemented in \textbf{RSiena}
as the weight matrix \(\mathbf{W}\) of the \texttt{XWX} effect);
component payoff functions are \(f_d\); and the logit precision
parameter is \(\beta > 0\).

A terminological note is warranted here, because the two ideas are
routinely conflated. The matrix \(\mathbf{E}\) is an \emph{influence
matrix} in the sense of \citet{RivkinSiggelkow2007}: it states which
components interact and with what sign, and it is an \emph{input}
supplied by the analyst. \emph{Epistasis} is the \emph{consequence} --
the dependence of one component's fitness contribution on the states of
others, which arises when a portfolio is evaluated through
\(\mathbf{E}\) as \(\mathbf{x}^{\top}\mathbf{E}\mathbf{x}\). The
distinction matters in \textbf{searchnet} because a third quantity sits
between the two. The component-component degree \(K_{CC}\) records how
far components are actually co-adopted at a given moment; it is computed
from the realized affiliation matrix and does not read \(\mathbf{E}\),
yet it evolves as a \emph{consequence} of \(\mathbf{E}\), which drives
the search that produces those affiliations. \(K_{CC}\) therefore
indexes the coupling epistasis acts through, generally differs from the
influence structure that generated it, and is not itself a measure of
the fitness consequence. Accordingly the package argument is named
\texttt{influence\_matrix}; \texttt{epistasis\_matrix} is retained as a
deprecated alias.

\subsection{Bipartite network formulation}\label{sec:bipartite}

The classical NK model represents a system configuration as an
\(N\)-dimensional binary vector \(\mathbf{x} = (x_1, x_2, \ldots, x_N)\)
where \(x_j \in \{0, 1\}\). \textbf{searchnet} generalizes this by
introducing \(M\) actors, each of whom maintains an affiliation vector.
The resulting system state is an \(M \times N\) binary matrix
\(\mathbf{B}\), where \(b_{ij} = 1\) indicates that actor \(i\)
affiliates with component \(j\). This matrix is the incidence matrix of
a bipartite network with actor nodes on one side and component nodes on
the other.

The bipartite representation has three immediate consequences. First, it
naturally accommodates \textbf{multiple agents} searching
simultaneously: each row of \(\mathbf{B}\) is an actor's configuration,
and actors' configurations interact through shared columns. Second, it
makes the \textbf{search environment endogenous}: the pattern of
affiliations determines which components are coupled (via shared actors)
and which actors are competitors (via shared components). Third, it
provides the structural substrate for the SAOM engine: \(\mathbf{B}\)
becomes the dependent variable in an evolving bipartite network model.

\subsection{\texorpdfstring{The \(\{K\}\) interpretive
framework}{The \textbackslash\{K\textbackslash\} interpretive framework}}\label{sec:k4}

The bipartite incidence matrix \(\mathbf{B}\) generates four coupled
degree statistics, each of which is individually a standard network
measure but whose joint dynamics expose the mechanism through which
micro-level search decisions produce macro-level landscape structure.
Table\textasciitilde{}\ref{tab:k4} summarizes the four dimensions.

\begin{table}[t]
\centering
\caption{The $\{K\}$ coupled degree framework.  Each dimension is a standard
network statistic; the contribution is their joint treatment as a coupled
system.}
\label{tab:k4}
\begin{tabular}{llll}
\toprule
Dimension & Symbol & Interpretation & Network measure \\
\midrule
Actor scope & $K_{AC}$ & Components per actor & Bipartite row degree \\
Component popularity & $K_{CA}$ & Actors per component & Bipartite column degree \\
Actor sociality & $K_{AA}$ & Co-affiliating actors per actor & Actor projection degree \\
Component epistasis & $K_{CC}$ & Co-occurring components per component & Component projection degree \\
\bottomrule
\end{tabular}
\end{table}

\begin{figure}
\centering
\includegraphics[width=0.8\linewidth,height=\textheight,keepaspectratio,alt={Empirical validation: K\_\{AA\} (competitive overlap) among U.S.~domestic airlines, 1993--2025 (Downing 2026a). The 4.8-fold increase in pairwise overlap demonstrates the \textbackslash\{K\textbackslash\} system's capacity to track competitive convergence dynamics using real network panel data.}]{searchnet-jss-NEW2_files/figures/cd2026_kaa_trajectory.png}
\caption{Empirical validation: \(K_{AA}\) (competitive overlap) among
U.S.~domestic airlines, 1993--2025 (Downing 2026a). The 4.8-fold
increase in pairwise overlap demonstrates the \(\{K\}\) system's
capacity to track competitive convergence dynamics using real network
panel data.}
\end{figure}

The classical NK parameter \(K\) maps directly onto \(K_{CC}\)---the
degree of a node in the component--component projection network. Two
components share a projection edge whenever at least one actor
affiliates with both, and the number of such edges incident on a focal
component determines how many other components epistatically affect its
fitness contribution. In the classical model, this coupling structure is
set exogenously by the modeler; in \textbf{searchnet}, it emerges
endogenously from the bipartite network and co-evolves with the other
three degree dimensions.

The analytical leverage of the \(\{K\}\) framework comes from treating
these four measures as a coupled system rather than examining any one in
isolation. Because all four derive from the same incidence matrix
\(\mathbf{B}\), they are structurally linked. A change in actor scope
(\(K_{AC}\)) mechanically alters component popularity (\(K_{CA}\)),
which in turn reshapes both the actor co-membership network (\(K_{AA}\))
and the component co-occurrence network (\(K_{CC}\)). The projections
are computed as:

\begin{align}
\mathbf{P}_{AA} &= \mathbf{B} \mathbf{B}^\top - \text{diag}(\mathbf{B} \mathbf{B}^\top), \label{eq:proj-actor} \\
\mathbf{P}_{CC} &= \mathbf{B}^\top \mathbf{B} - \text{diag}(\mathbf{B}^\top \mathbf{B}). \label{eq:proj-component}
\end{align}

The degree of actor \(i\) in the actor projection \(\mathbf{P}_{AA}\) is
\(K_{AA,i}
= \sum_{h \neq i} \mathbf{1}[(\mathbf{B}\mathbf{B}^\top)_{ih} > 0]\),
and symmetrically for \(K_{CC,j}\) in the component projection.
Analyzing any single measure in isolation obscures the cross-domain
dependencies that link actor behavior to landscape topology.

It is important to be precise about what the \(\{K\}\) framework
contributes. The four statistics are individually well known in the
bipartite network literature \citep{Latapy2008, BorgattiEverett1997}.
The contribution is not statistical but \textbf{interpretive}: \(\{K\}\)
provides a theoretical bridge between the NK landscape tradition in
strategy---where \(K\) is a scalar complexity parameter---and the
bipartite network tradition in sociology---where degree distributions
are descriptive summaries of affiliation structure. By grouping these
four measures and tracking their joint evolution within a SAOM, the
framework makes visible the transmission mechanism through which actors'
strategic choices propagate into landscape ruggedness.

\subsection{SAOM objective function specification}\label{sec:objective}

The SAOM objective function in \textbf{searchnet} governs actors'
affiliation decisions. At each micro-step, the selected actor \(i\)
evaluates the change in objective function value for each possible tie
change (creating or dissolving an affiliation with one component) and
selects the change that maximizes expected utility, subject to a
stochastic error term drawn from the type I extreme value distribution.

The objective function takes the form of
Equation\textasciitilde{}\ref{eq:saom-objective}, with the following key
statistics:

\begin{table}[t]
\centering
\caption{SAOM objective function statistics in \textbf{searchnet}.}
\label{tab:effects}
\begin{tabular}{llll}
\toprule
Statistic & RSiena name & Formula & Interpretation \\
\midrule
Density & \texttt{density} & $\sum_j b_{ij}$ & Baseline affiliation rate \\
Popularity & \texttt{inPop} & $\sum_j b_{ij} \sqrt{\sum_{h \neq i} b_{hj}}$ & Preference for popular components \\
Scope & \texttt{outAct} & $\sum_j b_{ij} \sqrt{\sum_{k \neq j} b_{ik}}$ & Scope expansion tendency \\
Epistasis & \texttt{XWX} & $\sum_j b_{ij} \sum_k w_{jk} b_{ik}$ & Exogenous component coupling \\
\bottomrule
\end{tabular}
\end{table}

The density effect serves as an intercept: a negative density parameter
creates baseline sparsity pressure, preventing actors from affiliating
with all components. The popularity effect (\texttt{inPop}) creates
preferential attachment to components that many other actors already
use. The scope effect (\texttt{outAct}) governs actors' tendency to
expand or contract the breadth of their affiliations. The epistasis
effect (\texttt{XWX}) introduces exogenous component--component coupling
via a modeler-supplied weight matrix \(\mathbf{W}\) (denoted
\(\mathbf{E}\) in the formal proofs; see
Section\textasciitilde{}\ref{sec:formal-properties}); actors derive
utility from affiliating with components that are coupled to their
existing affiliations.

Actor-level heterogeneity enters through constant covariates attached to
the objective function. A covariate vector \(\mathbf{v}\) (e.g.,
encoding strategic orientation) can modify the weight an actor places on
popularity seeking (\texttt{inPopX}), scope expansion (\texttt{egoX}),
or component characteristics (\texttt{altX}). This mechanism allows the
modeler to assign distinct strategic profiles to different actors---for
example, ``exploiters'' who preferentially attach to popular components
and ``explorers'' who expand into uncrowded niches.

The simulation proceeds via \textbf{RSiena}'s continuous-time Markov
chain algorithm. Actors are selected according to a rate function (by
default, equal rates for all actors), and the selected actor evaluates
all possible single-tie changes, choosing according to a multinomial
logit rule:

\begin{equation}\label{eq:choice}
  P(b_{ij} \to 1 - b_{ij}) =
  \frac{\exp\bigl(f_i(\mathbf{B}^{ij})\bigr)}
       {\sum_{j'} \exp\bigl(f_i(\mathbf{B}^{ij'})\bigr)},
\end{equation}

\noindent where \(\mathbf{B}^{ij}\) denotes the network configuration
obtained by toggling actor \(i\)'s tie to component \(j\). The
\texttt{search\_rsiena} method passes the full structure model to
\textbf{RSiena}'s \texttt{siena07} function, which executes the
continuous-time simulation internally. The
\texttt{iterations\_per\_actor} parameter controls the total number of
micro-steps: \(M \times\) \texttt{iterations\_per\_actor}.

For exogenous shocks, \textbf{searchnet} partitions the simulation
timeline into proportional segments. At each segment boundary, the
\(\theta\) parameters are updated to new values specified by the shock
schedule. The simulation then continues from the current network state
with the modified objective function, preserving path dependence while
allowing the researcher to model environmental discontinuities. This is
equivalent to running sequential \textbf{RSiena} simulations where each
wave's initial state is the terminal state of the previous wave.

\subsection{Mapping RSiena effects to strategic search
concepts}\label{sec:glossary}

The preceding objective function uses \textbf{RSiena} effect names that
are standard in social network analysis but may be unfamiliar to
strategy researchers. Table\textasciitilde{}\ref{tab:glossary} provides
a systematic mapping between RSiena effects available in
\textbf{searchnet} and their strategic search interpretations. Each
effect's parameter symbol links it to the corresponding \(\{K\}\)
dimension introduced in Section\textasciitilde{}\ref{sec:k4}, enabling
researchers to translate between the network estimation vocabulary and
the strategy-theoretic language of scope, complementarity, and
imitation.

\begin{table}[t]
\centering
\caption{Mapping RSiena effects to strategic search concepts.  Parameters correspond to the SAOM objective function (Equation~\ref{eq:saom-objective}).  The $\{K\}$ column indicates which coupled degree dimension each effect primarily governs.}
\label{tab:glossary}
\begin{tabular}{lllll}
\toprule
RSiena Effect & Parameter & $\{K\}$ Dim. & Strategy Concept & Interpretation \\
\midrule
\texttt{density}      & $\theta_{\text{dens}}$  & ---        & Scope cost               & Negative = costly to hold activities; controls portfolio breadth \\
\texttt{outAct}       & $\theta_{\text{scope}}$ & $K_{AC}$   & Diminishing returns       & Quadratic penalty on actor degree (scope) \\
\texttt{inPop}        & $\theta_{\text{pop}}$   & $K_{CA}$   & Bandwagon / legitimacy    & Positive = prefer popular components; preferential attachment \\
\texttt{XWX}          & $\theta_{\text{epist}}$ & $K_{CC}$   & Complementarity / synergy & Reward for holding epistatically linked components \\
\texttt{simEgoInDist2}& $\theta_{\text{imit}}$  & $K_{AA}$   & Imitation / herding       & Positive = actors converge to similar portfolios \\
\texttt{egoX}         & $\theta_{\text{strat}}$ & ---        & Strategy heterogeneity    & Actor-level covariate modifying search behavior \\
\texttt{altX}         & $\theta_{\text{value}}$ & ---        & Component value           & Exogenous activity-level payoff differences \\
\texttt{cycle4}       & $\theta_{\text{close}}$ & $K_{AA}$   & Structural closure        & Triadic closure in the projected social network \\
\bottomrule
\end{tabular}
\end{table}

The table makes explicit the division of labor among effects. The first
four rows---\texttt{density}, \texttt{outAct}, \texttt{inPop}, and
\texttt{XWX}---are the workhorses of most \textbf{searchnet}
specifications. Together they endogenize the four \(\{K\}\) dimensions:
scope cost sets the baseline, \texttt{outAct} introduces diminishing
returns to breadth, \texttt{inPop} creates demand-side herding toward
popular components, and \texttt{XWX} rewards configurations that exploit
exogenous complementarities. The remaining effects allow the researcher
to layer on additional mechanisms: imitation (\texttt{simEgoInDist2})
induces portfolio convergence among structurally similar actors, actor
and component covariates (\texttt{egoX}, \texttt{altX}) introduce
heterogeneity, and \texttt{cycle4} captures closure in the actor--actor
projection of the bipartite network. By selecting which effects to
include and setting their parameter values, the researcher controls
which strategic mechanisms are active in a given simulation---a design
choice whose consequences can be audited with the diagnostic tools
described in Section\textasciitilde{}\ref{sec:validation}.

\subsubsection{Formal properties}\label{sec:formal-properties}

The SaoMNK framework rests on formal results that establish its
relationship to the classical NK model and to quantal response
equilibrium. We state the main results here; complete proofs with all
lemmas are provided in the Online Appendix
\citep[see also][]{downing2026architecture} (available as
\texttt{inst/proofs/} in the package; access via
\texttt{searchnet\_proof()}).

\begin{description}

\item[Theorem 1 (Reduction).] The standard NK model is a degenerate special case
of the SAOM-NK model.  Specifically, the NK landscape and its adaptive walk are
recovered exactly under five restrictions: (R1)~$M = 1$ (single actor);
(R2)~$\mathbf{E} = \mathbf{A}$ (SAOM-NK influence matrix equals the NK
influence matrix); (R3)~$\boldsymbol{\theta} = \mathbf{0}$ (all SAOM
parameters zero); (R4)~$\beta \to \infty$ (deterministic greedy choice);
(R5)~uniform component selection probability.

\item[Theorem 2 (Generalization).] SAOM-NK strictly generalizes the NK model
along three independent dimensions: (a)~multi-actor endogenous interaction
($M > 1$), (b)~endogenous landscape co-evolution through SAOM network
statistics, and (c)~bounded rationality with tunable search precision
($\beta \in (0, \infty)$).  The coupled degree vector
$\{K\}= (K_{AC}, K_{CA}, K_{AA}, K_{CC})$ spans a strictly larger class of
fitness landscapes than the scalar $K$ parameterization.  The classical NK
model is recovered as a proper special case: NK $\subsetneq$ SAOM-NK.

\item[Theorem 3 (Approximation).] For any NK landscape
$(N, K, \mathbf{A}, \{f_d\}_{d=1}^N)$, there exists a parameterization
of the SAOM-NK model that reproduces the NK payoff structure to arbitrary
precision.  Three independent proof strategies establish this: (a)~constructive
existence via dummy covariates, (b)~actor heterogeneity substitution, and
(c)~universal approximation via mixed logit.

\item[Theorem 4 (SAOM--QRE Equivalence).] The SAOM multinomial logit choice
rule (Equation~\ref{eq:choice}) is isomorphic to a quantal response
equilibrium (QRE) in the bipartite affiliation game, establishing a
game-theoretic foundation for the statistical model.

\end{description}

\noindent Theorems\textasciitilde1--3 correspond to
Theorems\textasciitilde1--3 in \citet{downing2026architecture};
Theorem\textasciitilde4 is specific to this paper. These results are
applied empirically to U.S.~airline competitive dynamics in
\citet{downing2026architecture}.

\subsection{Why SAOM over ERGM?}\label{sec:saom-ergm}

A natural question is why \textbf{searchnet} adopts the SAOM rather than
the exponential random graph model
\citep[ERGM;][]{LusherKoskinenRobins2013} or its temporal extension
\citep[TERGM;][]{Hanneke2010} as its network engine. The choice follows
directly from the theoretical premise of the NK tradition.

The NK model is fundamentally about \textbf{agents making strategic
decisions}: an actor evaluates alternative configurations and selects
the one that maximizes a fitness criterion. The SAOM's rate-function
architecture is the natural formalization of this process. At each
micro-step, a single actor is selected, evaluates the marginal utility
of each possible tie change, and chooses probabilistically according to
a multinomial logit rule. This actor-centric, sequential-decision
structure maps directly onto the local search process in NK models,
where an agent evaluates neighboring configurations and moves to the
best alternative.

ERGMs, by contrast, characterize the equilibrium distribution of network
configurations without specifying a generative decision process. An ERGM
answers the question ``what network structures are over-represented
relative to chance?'' but does not model \emph{who} makes \emph{which}
decision \emph{when}. For a framework whose purpose is to endogenize
landscape ruggedness through actors' strategic choices, this equilibrium
focus omits precisely the mechanism of interest.

Beyond theoretical alignment, SAOMs offer practical advantages for the
simulation-based workflow that \textbf{searchnet} supports. The SAOM's
micro-step chain produces a complete trajectory of network states,
enabling the researcher to track how \(\{K\}\) dimensions evolve over
time, to identify critical transitions, and to introduce shocks at
precise points in the simulation. ERGM simulation, by contrast,
generates draws from a stationary distribution without temporal
ordering, making path-dependent analysis infeasible.

Finally, the SAOM's continuous-time formulation handles bipartite
networks with actor-level covariates naturally through \textbf{RSiena}'s
existing infrastructure. While bipartite ERGMs exist
\citep{WangPattison2009}, the SAOM's treatment of bipartite networks as
actor-driven affiliation processes is better suited to the
search-theoretic interpretation that motivates \textbf{searchnet}.

\section{\texorpdfstring{The \textbf{searchnet}
package}{The  package}}\label{sec:package}

\subsection{Architecture}\label{sec:architecture}

\textbf{searchnet} is built on \textsf{R6} classes organized in a
two-level inheritance hierarchy:

\begin{description}
\item[\texttt{SaomNkRSienaBiEnv\_base}] Core environment initialization,
  bipartite matrix management, graph construction, projection computation, and
  utility helpers (toggle, Jaccard, existence checks).  This base class
  manages the $M \times N$ incidence matrix, constructs the \textbf{RSiena} data
  objects, and provides low-level network operations.

\item[\texttt{SaomNkRSienaBiEnv}] Full simulation engine inheriting from the
  base class, adding the \texttt{search\_rsiena()} simulation method, chain
  statistics processing, $\{K\}$-4 degree computation, fitness landscape
  evaluation, shock handling, multi-wave experiment infrastructure, and the
  visualization interface.
\end{description}

Supporting modules provide separation of concerns:

\begin{itemize}
\item \texttt{saomnk-api.R}: User-facing functional API (\texttt{saomnk\_env},
  \texttt{saomnk\_model}, \texttt{saomnk\_run}) that wraps the R6 internals.
\item \texttt{nk-classic.R}: A self-contained implementation of the classical
  NK model, used as an external reference for the reduction check of
  Section~\ref{sec:validation}.
\item \texttt{saomnk-experiments.R}: \texttt{SaoMNKexperiments} R6 class for batch
  Monte Carlo simulation management.
\item \texttt{plot-*.R}: Eight modular plotting files (\texttt{plot-utility.R},
  \texttt{plot-degrees.R}, \texttt{plot-exploration.R}, \texttt{plot-markets.R},
  \texttt{plot-multiwave.R}, \texttt{plot-shocks.R}, \texttt{plot-snapshots.R},
  \texttt{plot-phase-space.R}) containing visualization functions.
\item \texttt{searchnet-game.R}: Strategy game mode providing an interactive
  game loop where a human player controls one firm and competes against AI
  opponents that follow SAOM-NK logit choice rules.
\item \texttt{searchnet-teaching.R}: Classroom teaching module for Capsim-style
  strategy simulations supporting multiple industry presets, difficulty
  levels, and exogenous shocks.
\item \texttt{utils.R}: Standalone helper functions.
\end{itemize}

The package follows standard \textsf{R} packaging conventions with
\texttt{DESCRIPTION}, \texttt{NAMESPACE}, a \textbf{testthat} test suite
(300 tests across 23 files), and eight vignettes. All exported functions
are prefixed with \texttt{saomnk\_} or \texttt{searchnet\_} to avoid
namespace collisions.

\subsection{User API}\label{sec:api}

A central design goal of \textbf{searchnet} is to provide a clean
functional interface that shields users from the complexity of
\textbf{RSiena}'s internal naming conventions and the R6 class
hierarchy. The workflow follows three steps: environment creation, model
specification, and simulation execution.

\subsubsection{Step 1: Create the search
environment}\label{step-1-create-the-search-environment}

The \texttt{saomnk\_env()} function initializes a bipartite search
environment:

\begin{Shaded}
\begin{Highlighting}[]
\NormalTok{env }\OtherTok{\textless{}{-}} \FunctionTok{saomnk\_env}\NormalTok{(}\AttributeTok{M =} \DecValTok{8}\NormalTok{, }\AttributeTok{N =} \DecValTok{12}\NormalTok{, }\AttributeTok{density =} \DecValTok{0}\NormalTok{, }\AttributeTok{seed =} \DecValTok{42}\NormalTok{)}
\end{Highlighting}
\end{Shaded}

This creates a \texttt{SaomNkRSienaBiEnv} object representing 8 actors
and 12 components, starting from an empty network (all actors
unaffiliated). The \texttt{seed} argument ensures reproducibility of the
initial network draw.

\subsubsection{Step 2: Specify the structure
model}\label{step-2-specify-the-structure-model}

The \texttt{saomnk\_model()} function builds the SAOM objective function
specification using intuitive parameter names:

\begin{Shaded}
\begin{Highlighting}[]
\DocumentationTok{\#\# Define a modular influence matrix: 4 blocks of 3 components}
\NormalTok{K\_matrix }\OtherTok{\textless{}{-}} \FunctionTok{saomnk\_block\_diagonal}\NormalTok{(}\DecValTok{12}\NormalTok{, }\DecValTok{4}\NormalTok{)}

\DocumentationTok{\#\# Specify the model}
\NormalTok{mod }\OtherTok{\textless{}{-}} \FunctionTok{saomnk\_model}\NormalTok{(}
  \AttributeTok{density    =} \SpecialCharTok{{-}}\FloatTok{0.5}\NormalTok{,     }\CommentTok{\# sparsity pressure}
  \AttributeTok{popularity =}  \FloatTok{0.2}\NormalTok{,     }\CommentTok{\# preferential attachment}
  \AttributeTok{scope      =}  \FloatTok{0.1}\NormalTok{,     }\CommentTok{\# scope expansion}
  \AttributeTok{influence\_matrix =}\NormalTok{ K\_matrix,}
  \AttributeTok{influence\_weight =} \FloatTok{0.05}
\NormalTok{)}
\NormalTok{mod}
\CommentTok{\#\textgreater{} SaoMNK Structure Model}
\CommentTok{\#\textgreater{} {-}{-}{-}{-}{-}{-}{-}{-}{-}{-}{-}{-}{-}{-}{-}{-}{-}{-}{-}{-}{-}{-}}
\CommentTok{\#\textgreater{} Effects on the bipartite network DV (3):}
\CommentTok{\#\textgreater{}   density (density)      theta =   {-}0.5  [fixed]}
\CommentTok{\#\textgreater{}   popularity (inPop)     theta =    0.2  [fixed]}
\CommentTok{\#\textgreater{}   scope (outAct)         theta =    0.1  [fixed]}
\CommentTok{\#\textgreater{} Influence / dyadic covariates (1):}
\CommentTok{\#\textgreater{}   epistasis (XWX)        theta =   0.05  [fixed]  [influence matrix W, 12 x 12]}
\end{Highlighting}
\end{Shaded}

Users never need to write internal references such as
\texttt{"self\$bipartite\_rsienaDV"} or construct nested covariate lists
manually. The function assembles the complete \textbf{RSiena} structure
model internally. Actor strategy heterogeneity is specified through the
\texttt{strategies} argument:

\begin{Shaded}
\begin{Highlighting}[]
\DocumentationTok{\#\# Heterogeneous actors: varying scope and popularity preferences}
\NormalTok{mod\_strat }\OtherTok{\textless{}{-}} \FunctionTok{saomnk\_model}\NormalTok{(}
  \AttributeTok{density    =} \SpecialCharTok{{-}}\FloatTok{0.3}\NormalTok{,}
  \AttributeTok{popularity =}  \FloatTok{0.2}\NormalTok{,}
  \AttributeTok{scope      =}  \FloatTok{0.1}\NormalTok{,}
  \AttributeTok{influence\_matrix =}\NormalTok{ K\_matrix,}
  \AttributeTok{influence\_weight =} \FloatTok{0.05}\NormalTok{,}
  \AttributeTok{strategies =} \FunctionTok{list}\NormalTok{(}
    \AttributeTok{egoX   =} \FunctionTok{rep}\NormalTok{(}\FunctionTok{c}\NormalTok{(}\SpecialCharTok{{-}}\DecValTok{1}\NormalTok{, }\DecValTok{0}\NormalTok{, }\DecValTok{1}\NormalTok{, }\DecValTok{0}\NormalTok{), }\AttributeTok{length.out =} \DecValTok{8}\NormalTok{),}
    \AttributeTok{inPopX =} \FunctionTok{rep}\NormalTok{(}\FunctionTok{c}\NormalTok{( }\DecValTok{1}\NormalTok{, }\DecValTok{0}\NormalTok{,}\SpecialCharTok{{-}}\DecValTok{1}\NormalTok{, }\DecValTok{0}\NormalTok{), }\AttributeTok{length.out =} \DecValTok{8}\NormalTok{)}
\NormalTok{  )}
\NormalTok{)}
\end{Highlighting}
\end{Shaded}

\subsubsection{Step 3: Run the
simulation}\label{step-3-run-the-simulation}

The \texttt{saomnk\_run()} function executes the SAOM simulation:

\begin{Shaded}
\begin{Highlighting}[]
\FunctionTok{saomnk\_run}\NormalTok{(env, mod, }\AttributeTok{steps\_per\_actor =} \DecValTok{30}\NormalTok{, }\AttributeTok{seed =} \DecValTok{12345}\NormalTok{)}
\end{Highlighting}
\end{Shaded}

The function passes the structure model to the internal
\texttt{search\_rsiena()} method, runs the continuous-time SAOM
simulation via \textbf{RSiena}, extracts chain statistics, and computes
the \(\{K\}\)-4 degree trajectories. The \texttt{env} object is modified
in place (standard R6 semantics) and returned invisibly.

\subsection{Visualization}\label{sec:visualization}

\textbf{searchnet} provides a comprehensive visualization toolkit
organized around four analytical perspectives.

\subsubsection{\texorpdfstring{The \(\{K\}\)-4
panel}{The \textbackslash\{K\textbackslash\}-4 panel}}\label{the-k-4-panel}

The signature visualization is the coupled degree panel, which displays
all four \(\{K\}\) dimensions over the simulation timeline:

\begin{Shaded}
\begin{Highlighting}[]
\FunctionTok{saomnk\_plot\_k4}\NormalTok{(env, }\AttributeTok{smooth =} \FloatTok{0.3}\NormalTok{)}
\end{Highlighting}
\end{Shaded}

\begin{figure}

{\centering \includegraphics{searchnet-jss-NEW2_files/figure-latex/k4-panel-1} 

}

\caption{The $\{K\}$-4 coupled degree panel showing actor scope ($K_{AC}$), component popularity ($K_{CA}$), actor sociality ($K_{AA}$), and component epistasis ($K_{CC}$) trajectories over 240 simulation micro-steps.  Each line represents a single actor or component; the heavy line is the LOESS-smoothed mean.}\label{fig:k4-panel}
\end{figure}

The panel reveals the structural coupling between dimensions: as actors
expand scope (\(K_{AC}\)), they mechanically increase component
popularity (\(K_{CA}\)), which in turn elevates both actor sociality
(\(K_{AA}\)) and component epistasis (\(K_{CC}\)). This coupling is the
core mechanism through which individual search decisions endogenize the
fitness landscape.

\includegraphics[width=0.45\linewidth,height=\textheight,keepaspectratio,alt={Application breadth: (left) Startup workforce \textbackslash\{K\textbackslash\}-4 panel showing scope, popularity, sociality, and epistasis dynamics for 8 startups across 12 skill categories; (right) Phase space trajectory showing startup evolution through Network \textbackslash times Behavior \textbackslash times Fitness space.}]{searchnet-jss-NEW2_files/figures/bort_k4_panel.png}
\includegraphics[width=0.45\linewidth,height=\textheight,keepaspectratio]{searchnet-jss-NEW2_files/figures/bort_phase_space.png}

\subsubsection{Network snapshots}\label{network-snapshots}

The \texttt{saomnk\_plot\_snapshots()} function renders the bipartite
network at selected simulation steps, along with its unipartite
projections:

\begin{Shaded}
\begin{Highlighting}[]
\FunctionTok{saomnk\_plot\_snapshots}\NormalTok{(env, }\AttributeTok{steps =} \FunctionTok{c}\NormalTok{(}\DecValTok{1}\NormalTok{, }\DecValTok{120}\NormalTok{, }\DecValTok{240}\NormalTok{))}
\end{Highlighting}
\end{Shaded}

\begin{figure}

{\centering \includegraphics{searchnet-jss-NEW2_files/figure-latex/snapshots-1} 

}

\caption{Network snapshots at the initial state, midpoint, and final state of the simulation.  Top row: bipartite actor-component network.  Bottom rows: actor co-membership projection ($K_{AA}$) and component co-occurrence projection ($K_{CC}$).}\label{fig:snapshots-1}
\end{figure}
\begin{figure}

{\centering \includegraphics{searchnet-jss-NEW2_files/figure-latex/snapshots-2} 

}

\caption{Network snapshots at the initial state, midpoint, and final state of the simulation.  Top row: bipartite actor-component network.  Bottom rows: actor co-membership projection ($K_{AA}$) and component co-occurrence projection ($K_{CC}$).}\label{fig:snapshots-2}
\end{figure}
\begin{figure}

{\centering \includegraphics{searchnet-jss-NEW2_files/figure-latex/snapshots-3} 

}

\caption{Network snapshots at the initial state, midpoint, and final state of the simulation.  Top row: bipartite actor-component network.  Bottom rows: actor co-membership projection ($K_{AA}$) and component co-occurrence projection ($K_{CC}$).}\label{fig:snapshots-3}
\end{figure}
\begin{figure}

{\centering \includegraphics{searchnet-jss-NEW2_files/figure-latex/snapshots-4} 

}

\caption{Network snapshots at the initial state, midpoint, and final state of the simulation.  Top row: bipartite actor-component network.  Bottom rows: actor co-membership projection ($K_{AA}$) and component co-occurrence projection ($K_{CC}$).}\label{fig:snapshots-4}
\end{figure}

\subsubsection{Utility decomposition}\label{utility-decomposition}

The \texttt{saomnk\_plot\_utility()} function decomposes each actor's
total objective function value into contributions from individual
network statistics:

\begin{Shaded}
\begin{Highlighting}[]
\FunctionTok{saomnk\_plot\_utility}\NormalTok{(env, }\AttributeTok{smooth =} \FloatTok{0.35}\NormalTok{)}
\end{Highlighting}
\end{Shaded}

\begin{figure}

{\centering \includegraphics{searchnet-jss-NEW2_files/figure-latex/utility-1} 

}

\caption{Utility decomposition showing each actor's objective function value broken into contributions from density, popularity, scope, and epistasis effects over the simulation timeline.}\label{fig:utility}
\end{figure}

This decomposition provides direct insight into the strategic trade-offs
governing search behavior: an actor whose utility is dominated by the
epistasis term is searching for complementary components, while one
dominated by the popularity term is following the crowd.

\subsubsection{3D phase space}\label{d-phase-space}

The \texttt{saomnk\_plot\_phase\_space\_3d()} function renders the
simulation trajectory in a three-dimensional Network \(\times\) Behavior
\(\times\) Fitness space, where the three axes correspond to structural
properties (e.g., \(K_{AC}\) or density), behavioral properties (e.g.,
exploration rate or scope change), and outcome measures (utility or NK
fitness). The resulting ``SCP(t) trajectory'\,' traces each actor's path
through this space over simulation time, providing an integrated view of
how structural position, behavioral strategy, and performance co-evolve.
Supporting functions include \texttt{saomnk\_plot\_phase\_heatmap()} for
density-based projections, \texttt{saomnk\_plot\_phase\_evolution()} for
animated trajectories, \texttt{saomnk\_plot\_phase\_comparison()} for
juxtaposing multiple simulation runs, and
\texttt{searchnet\_export\_phase\_space()} for Manim-compatible data
export.

\subsection{Validation against the classical NK
model}\label{sec:validation}

A software paper should demonstrate that its engine reproduces the model
it generalises. \textbf{searchnet} ships a self-contained implementation
of the classical Kauffman NK model, independent of \textbf{RSiena} and
of the SaoMNK engine, and uses it as an external reference.

\subsubsection{A conventional NK
implementation}\label{a-conventional-nk-implementation}

\texttt{nk\_landscape()} constructs a Kauffman landscape by exhaustive
enumeration of all \(2^N\) configurations. Each locus \(i\) contributes
\(f_i(b_i,
b_{\Omega(i)})\), drawn i.i.d.~\(U(0,1)\) for every distinct pattern of
its epistatic neighbourhood \(\Omega(i)\); total fitness is the mean of
the \(N\) contributions. Three neighbourhood topologies are supported:
\texttt{"adjacent"} (Kauffman's ring model), \texttt{"random"}, and
\texttt{"block"} (near-decomposable). Companion functions provide
adaptive walks (\texttt{nk\_walk()}, with steepest-ascent,
random-ascent, and random-improving rules), exhaustive local optimum
enumeration (\texttt{nk\_local\_optima()}), and the canonical ruggedness
sweep (\texttt{nk\_sweep\_K()}).

These functions are useful in their own right, but their role here is to
supply a reference implementation that shares no code with the SAOM
engine.

\subsubsection{Verifying the reduction}\label{verifying-the-reduction}

Theorem 1 states that the classical NK model is the \(M = 1\),
\(\beta \rightarrow \infty\) special case of SaoMNK.
\texttt{nk\_to\_saomnk()} emits the environment and structure-model
settings that place a single-actor SaoMNK environment on a supplied NK
landscape, making the reduction constructive rather than merely
asserted.

\texttt{nk\_verify\_reduction()} then checks the construction. It
exploits the defining property of the NK model: contribution \(f_i\)
depends only on the loci in \(\Omega(i)\), so any two configurations
sharing the same neighbourhood pattern must carry identical \(f_i\).
Rebuilding each locus's payoff table keyed by that pattern and
recomputing fitness must return the original values exactly. A non-zero
difference would indicate an indexing or bit-ordering error in landscape
construction.

\begin{Shaded}
\begin{Highlighting}[]
\NormalTok{nk }\OtherTok{\textless{}{-}} \FunctionTok{nk\_landscape}\NormalTok{(}\AttributeTok{N =} \DecValTok{10}\NormalTok{, }\AttributeTok{K =} \DecValTok{3}\NormalTok{, }\AttributeTok{seed =} \DecValTok{42}\NormalTok{)}
\FunctionTok{nk\_verify\_reduction}\NormalTok{(}\AttributeTok{N =} \DecValTok{10}\NormalTok{, }\AttributeTok{K =} \DecValTok{3}\NormalTok{, }\AttributeTok{seed =} \DecValTok{42}\NormalTok{)}
\CommentTok{\#\textgreater{} NK reduction check: N=10, K=3, model="adjacent"}
\CommentTok{\#\textgreater{}   configurations   : 1,024}
\CommentTok{\#\textgreater{}   max |difference| : 0.00e+00}
\CommentTok{\#\textgreater{}   result           : PASS}
\end{Highlighting}
\end{Shaded}

The check passes at machine precision (maximum absolute difference
\(0.00 \times 10^{0}\)) across the \((N, K)\) combinations exercised in
the test suite. Configuration encoding is most-significant-bit-first
throughout and is covered by exhaustive round-trip tests for
\(N = 1, \ldots, 12\); this guards a class of bit-ordering error that is
easy to introduce and silent when it occurs.

\subsubsection{Recovering classical
results}\label{recovering-classical-results}

The implementation reproduces the standard qualitative findings of the
NK literature. With \(K = 0\) the landscape is single-peaked, and a
steepest-ascent walk from any starting configuration reaches the global
optimum. As \(K\) rises the number of local optima grows, approaching
the familiar \(2^N / (K+1)\) scaling, adaptive walks terminate sooner,
and terminal fitness falls
\citep{Kauffman1987, KauffmanWeinberger1989, Levinthal1997}.
\texttt{nk\_sweep\_K()} reproduces this experiment directly.

\subsection{Manim animation pipeline}\label{sec:manim}

For presentation and pedagogical purposes, \textbf{searchnet} exports
simulation data in formats compatible with the \textbf{Manim} animation
library \citep{Manim2022}. The exported data includes per-step bipartite
matrices, projection networks, degree trajectories, and utility
decompositions. The package ships with 16 ready-to-render Manim scene
classes organized across six Python modules: bipartite network
evolution, \(\{K\}\)-4 degree trajectories, 3D fitness landscapes, shock
response dynamics, 3D phase space trajectories, a seven-panel research
dashboard (covering \(\{K\}\) evolution, bipartite structure, fitness
surfaces, information funnels, shock response, and game-theoretic
equilibria), and four Blume--Gibbs scenes that animate the
statistical-mechanical foundations of convergence (Ising--logit
correspondence, temperature dials, Gibbs potential surfaces, and
convergence diagnostics). This pipeline is accessed through the
\texttt{saomnk\_export\_manim()} method on the environment object and
the \texttt{searchnet\_export\_phase\_space()} function for phase space
data.

\includegraphics[width=0.45\linewidth,height=\textheight,keepaspectratio,alt={Manim animation capabilities: (a) Basin geometry narrative showing 4 policy types diverging to different basins on a 3D NK landscape; (b) Convergence climb showing airline agents ascending a rugged fitness surface; (c) Blume Ising logit clustering; (d) ELO tournament rating evolution.}]{searchnet-jss-NEW2_files/figures/manim_basin_narrative.png}
\includegraphics[width=0.45\linewidth,height=\textheight,keepaspectratio]{searchnet-jss-NEW2_files/figures/manim_convergence_climb.png}

\subsection{Interactive features}\label{sec:interactive}

\textbf{searchnet} includes two interactive modules that extend the
simulation engine from a research tool to an instructional platform.

\textbf{Strategy game mode.} The \texttt{searchnet\_game\_*} family of
functions (five exported functions) provides an interactive game loop
where a human player controls one firm and competes against AI opponents
that follow the SAOM conditional-logit choice rule. The player
initializes a game with \texttt{searchnet\_game\_init()}, specifying the
number of firms, difficulty level (\texttt{"easy"}, \texttt{"medium"},
\texttt{"hard"}, mapping to logit precision
\(\beta \in \{0.5, 2, 10\}\)), and game objective (maximize fitness,
minimize competitive overlap, survive an exogenous shock, or outperform
the quantal response equilibrium prediction). At each round, the player
queries available moves via
\texttt{searchnet\_game\_available\_moves()}, selects an action, and
advances the simulation with \texttt{searchnet\_game\_step()}. AI
opponents respond simultaneously according to the SAOM logit rule, and
the scoreboard (\texttt{searchnet\_game\_scoreboard()}) tracks all
firms' \(\{K\}\) statistics. After the final round,
\texttt{searchnet\_game\_summary()} compiles trajectory data suitable
for replay visualization.

\textbf{Classroom teaching module.} The \texttt{searchnet\_classroom\_*}
family (six exported functions) supports Capsim-style strategy
simulations for classroom use. An instructor initializes a session with
\texttt{searchnet\_classroom\_init()}, choosing from industry presets
(airline, tech, pharma) or custom parameters, and specifying a
difficulty level that controls the number of AI firms and optional
exogenous shocks. Each student controls one firm, submitting portfolio
decisions (add or drop activities) through
\texttt{searchnet\_classroom\_submit()}. Once all students have
submitted, \texttt{searchnet\_classroom\_advance()} executes decisions
simultaneously, runs AI opponent responses via SAOM ministeps, and
updates the class leaderboard. At the end of the game,
\texttt{searchnet\_classroom\_debrief()} generates a suite of reports:
final rankings, round-by-round histories, \(\{K\}\)-4 panels, decision
logs, and shock records.

\subsection{Proof registry}\label{sec:proofs}

The package ships with a complete formal verification chain for the
SaoMNK framework, accessible via \texttt{searchnet\_proof()}. The proof
registry (\texttt{inst/proofs/PROOF\_TABLE.md}) enumerates 81 logical
steps---from axioms through definitions, lemmas, and propositions to the
four main theorems---organized so that every result can be traced
backward through its dependencies to foundational axioms. This
self-contained proof chain provides reviewers and users with a
verifiable logical scaffold for the theoretical claims summarized in
Section\textasciitilde{}\ref{sec:formal-properties}.

\section{Illustrations}\label{sec:illustrations}

\subsection{Reproducing Levinthal (1997)}\label{sec:levinthal}

The first illustration establishes that \textbf{searchnet} can reproduce
the core dynamics of the classical NK model under appropriate parameter
settings. We configure an environment with no endogenous network
effects, leaving only the exogenous influence matrix active. Setting
both \texttt{popularity} and \texttt{scope} to zero removes every term
through which one actor's state can enter another's objective function,
so the actors search \emph{independently}: the model reduces to parallel
single-searcher processes on a landscape whose ruggedness is determined
solely by the exogenous coupling matrix, which is the functional
equivalent of the classical NK search process. We use \(M = 2\) rather
than \(M = 1\) because \texttt{search\_rsiena()} requires at least two
actors --- \textbf{RSiena}'s \texttt{sienaDataCreate()} does not accept
a single-actor bipartite network. For genuinely single-actor work the
landscape methods (\texttt{compute\_fitness\_landscape},
\texttt{verify\_nk\_equivalence}) accept \(M = 1\) directly, and
Section\textasciitilde{} ef\{sec:validation\} uses them for exactly
that.

\begin{Shaded}
\begin{Highlighting}[]
\DocumentationTok{\#\# Two INDEPENDENT searchers (popularity = scope = 0 removes all coupling),}
\DocumentationTok{\#\# 12 components, empty start.  M \textgreater{}= 2 is required by RSiena\textquotesingle{}s bipartite}
\DocumentationTok{\#\# data constructor; with the interaction terms zeroed the actors do not}
\DocumentationTok{\#\# influence one another, so each is a classical NK searcher.}
\NormalTok{env\_nk }\OtherTok{\textless{}{-}} \FunctionTok{saomnk\_env}\NormalTok{(}\AttributeTok{M =} \DecValTok{2}\NormalTok{, }\AttributeTok{N =} \DecValTok{12}\NormalTok{, }\AttributeTok{density =} \DecValTok{0}\NormalTok{, }\AttributeTok{seed =} \DecValTok{1234}\NormalTok{)}

\DocumentationTok{\#\# Block{-}diagonal epistasis: 3 modules of 4 components (K \textasciitilde{} 3)}
\NormalTok{K\_levinthal }\OtherTok{\textless{}{-}} \FunctionTok{saomnk\_block\_diagonal}\NormalTok{(}\DecValTok{12}\NormalTok{, }\DecValTok{3}\NormalTok{)}

\DocumentationTok{\#\# Only epistasis active; no endogenous effects}
\NormalTok{mod\_nk }\OtherTok{\textless{}{-}} \FunctionTok{saomnk\_model}\NormalTok{(}
  \AttributeTok{density    =} \SpecialCharTok{{-}}\FloatTok{0.3}\NormalTok{,}
  \AttributeTok{popularity =} \DecValTok{0}\NormalTok{,}
  \AttributeTok{scope      =} \DecValTok{0}\NormalTok{,}
  \AttributeTok{influence\_matrix =}\NormalTok{ K\_levinthal,}
  \AttributeTok{influence\_weight =} \FloatTok{0.15}
\NormalTok{)}

\FunctionTok{saomnk\_run}\NormalTok{(env\_nk, mod\_nk, }\AttributeTok{steps\_per\_actor =} \DecValTok{100}\NormalTok{, }\AttributeTok{seed =} \DecValTok{42}\NormalTok{)}
\end{Highlighting}
\end{Shaded}

Under this parameterization, the agent's search trajectory exhibits the
signature features of NK dynamics: rapid initial improvement as the
agent discovers complementary component clusters, followed by
convergence to a local optimum where no single-component change improves
utility. The K-4 panel reduces to the \(K_{AC}\) trajectory (actor
scope) and the \(K_{CC}\) trajectory (epistasis), with the latter
reflecting the fixed block-diagonal structure; the \(K_{CA}\) and
\(K_{AA}\) channels are flat by construction, since nothing couples the
two searchers.

\begin{Shaded}
\begin{Highlighting}[]
\NormalTok{env\_nk}\SpecialCharTok{$}\FunctionTok{plot\_actor\_degrees}\NormalTok{(}\AttributeTok{loess\_span =} \FloatTok{0.5}\NormalTok{)}
\end{Highlighting}
\end{Shaded}

\begin{figure}

{\centering \includegraphics{searchnet-jss-NEW2_files/figure-latex/levinthal-degrees-1} 

}

\caption{Actor scope ($K_{AC}$) trajectory for independent agents searching on a block-diagonal epistasis landscape.  The trajectory shows rapid scope expansion followed by stabilization---the SAOM equivalent of convergence to a local optimum on a rugged NK landscape.}\label{fig:levinthal-degrees}
\end{figure}

This result provides face validity: when all endogenous network effects
are disabled, \textbf{searchnet} recovers the qualitative behavior of
the classical NK model. The landscape is effectively exogenous
(determined by the fixed \(\mathbf{W}\) matrix), search is local
(single-tie changes), and convergence reflects the trapping property of
rugged landscapes.

A deeper convergence analysis is provided in the Blume tutorial vignette
(\texttt{vignette("saomnk-blume-tutorial")}), which establishes the
statistical-mechanical foundations of the SAOM-NK choice rule by
connecting it to the Blume--Capel model \citep{Blume1993}, Gibbs
measures, and the Ising--logit correspondence.

\subsubsection{Ergodicity: forgetting the initial
condition}\label{sec:ergodicity}

Theorem\textasciitilde4 asserts that the chain's stationary distribution
is unique and independent of where the simulation starts. That property
is easy to illustrate badly. Running the model once from a sparse start
and once from a dense start, then observing that the two final densities
are close, establishes very little: no threshold was declared in
advance, a single run per arm supplies no sampling distribution,
and---most importantly---an ergodic chain forgets its initial condition
\emph{at a rate}, which a pair of endpoints discards.

\textbf{searchnet} therefore ships the demonstration as a measurement.
\texttt{searchnet\_ergodicity\_sweep()} runs the same fixed-coefficient
model from two contrasting starting densities across a range of chain
lengths, with independent replicates at each, and reports how fast the
between-arm gap decays:

\begin{Shaded}
\begin{Highlighting}[]
\NormalTok{erg }\OtherTok{\textless{}{-}} \FunctionTok{searchnet\_ergodicity\_sweep}\NormalTok{(}
  \AttributeTok{M =} \DecValTok{12}\NormalTok{, }\AttributeTok{N =} \DecValTok{15}\NormalTok{, }\AttributeTok{start\_densities =} \FunctionTok{c}\NormalTok{(}\FloatTok{0.1}\NormalTok{, }\FloatTok{0.8}\NormalTok{),}
  \AttributeTok{run\_lengths =} \FunctionTok{c}\NormalTok{(}\DecValTok{15}\NormalTok{, }\DecValTok{30}\NormalTok{, }\DecValTok{60}\NormalTok{, }\DecValTok{120}\NormalTok{, }\DecValTok{240}\NormalTok{),}
  \AttributeTok{replicates =} \DecValTok{10}\NormalTok{, }\AttributeTok{equivalence\_margin =} \FloatTok{0.05}\NormalTok{, }\AttributeTok{seed =} \DecValTok{42}
\NormalTok{)}
\NormalTok{erg}
\CommentTok{\#\textgreater{} Ergodicity sweep (Theorem 4: independence from initial conditions)}
\CommentTok{\#\textgreater{}   12 actors x 15 components (180 cells; one tie = 0.0056 density)}
\CommentTok{\#\textgreater{}   arms starting at 0.10 and 0.80, 10 replicates each}
\CommentTok{\#\textgreater{} }
\CommentTok{\#\textgreater{}   run length   mean(lo)   mean(hi)        gap}
\CommentTok{\#\textgreater{}           15     0.4367     0.7750     0.3383}
\CommentTok{\#\textgreater{}           30     0.5128     0.6906     0.1778}
\CommentTok{\#\textgreater{}           60     0.4761     0.5867     0.1106}
\CommentTok{\#\textgreater{}          120     0.4872     0.5150     0.0278}
\CommentTok{\#\textgreater{}          240     0.5072     0.5106     0.0033}
\CommentTok{\#\textgreater{} }
\CommentTok{\#\textgreater{}   initial gap 0.7000 {-}{-}\textgreater{} final gap 0.0033  (99.5\% collapse)}
\CommentTok{\#\textgreater{}   Monte Carlo floor (10 replicates): 0.0137}
\CommentTok{\#\textgreater{}   decay above the floor: log(gap) \textasciitilde{} {-}1.15 * log(iterations), R\^{}2 = 0.943}
\CommentTok{\#\textgreater{}   (fit on 4 point(s); 240 excluded as at{-}or{-}below the floor)}
\CommentTok{\#\textgreater{} }
\CommentTok{\#\textgreater{}   TOST at 240 iterations, margin +/{-} 0.050, 90\% CI [{-}0.0331, 0.0265]:}
\CommentTok{\#\textgreater{}   EQUIVALENT {-}{-} the arms are statistically indistinguishable at the}
\CommentTok{\#\textgreater{}   declared margin. The chain has forgotten where it started.}
\FunctionTok{plot}\NormalTok{(erg)}
\end{Highlighting}
\end{Shaded}

\begin{figure}

{\centering \includegraphics{searchnet-jss-NEW2_files/figure-latex/ergodicity-1} 

}

\caption{Independence from initial conditions, measured rather than asserted. (a) Two arms starting at density 0.10 and 0.80 converge as run length grows; points are replicates. (b) The between-arm gap decays as a power of run length on log--log axes. The dashed line is the declared equivalence margin; the shaded band is the Monte Carlo floor, below which the measured gap reflects replicate noise rather than mixing, and which is excluded from the fitted slope.}\label{fig:ergodicity}
\end{figure}

Two arms separated by a starting gap of 0.70 close to within 0.01 by 240
iterations per actor, a collapse of more than 99\%, and the gap decays
as a power of run length with slope near \(-1\) and \(R^2 > 0.9\). The
verdict is computed rather than printed: at the longest run length the
arms are compared by a two-one-sided-tests equivalence procedure against
a margin fixed in the call, so the function is able to report
\emph{non}-equivalence, and does so at the shorter run lengths. This is
what makes the sweep evidence rather than illustration.

Two details matter for anyone reproducing this. The Monte Carlo floor is
drawn explicitly, because once the arms are genuinely equivalent the
measured gap settles at the expected difference between two sample means
from the same distribution rather than falling to zero; fitting a decay
slope through those points would overstate the rate. And the grid is
sized so that density has resolution: on the 4 \(\times\) 5 grid used in
earlier versions of this demonstration a single tie moved the density by
0.05, which is coarser than any equivalence margin worth declaring.

The vignette also provides computational verification that the SAOM
Markov chain satisfies detailed balance with respect to the Gibbs
distribution, which is the mechanism underlying the convergence measured
here.

\subsection{Endogenous landscape with heterogeneous
strategies}\label{sec:endogenous}

The second illustration activates the full \textbf{searchnet} machinery
to demonstrate endogenous landscape formation. We configure 12 actors
searching over 12 components with active density, popularity, scope, and
epistasis effects, plus heterogeneous strategy covariates.

\begin{Shaded}
\begin{Highlighting}[]
\NormalTok{env\_endog }\OtherTok{\textless{}{-}} \FunctionTok{saomnk\_env}\NormalTok{(}\AttributeTok{M =} \DecValTok{12}\NormalTok{, }\AttributeTok{N =} \DecValTok{12}\NormalTok{, }\AttributeTok{density =} \DecValTok{0}\NormalTok{, }\AttributeTok{seed =} \DecValTok{1234}\NormalTok{)}

\NormalTok{K\_modular }\OtherTok{\textless{}{-}} \FunctionTok{saomnk\_block\_diagonal}\NormalTok{(}\DecValTok{12}\NormalTok{, }\DecValTok{4}\NormalTok{)}

\NormalTok{mod\_endog }\OtherTok{\textless{}{-}} \FunctionTok{saomnk\_model}\NormalTok{(}
  \AttributeTok{density    =} \SpecialCharTok{{-}}\FloatTok{0.3}\NormalTok{,}
  \AttributeTok{popularity =}  \FloatTok{0.2}\NormalTok{,}
  \AttributeTok{scope      =}  \FloatTok{0.1}\NormalTok{,}
  \AttributeTok{influence\_matrix =}\NormalTok{ K\_modular,}
  \AttributeTok{influence\_weight =} \FloatTok{0.02}\NormalTok{,}
  \AttributeTok{strategies =} \FunctionTok{list}\NormalTok{(}
    \AttributeTok{egoX   =} \FunctionTok{rep}\NormalTok{(}\FunctionTok{c}\NormalTok{(}\SpecialCharTok{{-}}\DecValTok{1}\NormalTok{, }\DecValTok{0}\NormalTok{, }\DecValTok{1}\NormalTok{), }\AttributeTok{length.out =} \DecValTok{12}\NormalTok{),}
    \AttributeTok{inPopX =} \FunctionTok{rep}\NormalTok{(}\FunctionTok{c}\NormalTok{( }\DecValTok{1}\NormalTok{, }\DecValTok{0}\NormalTok{,}\SpecialCharTok{{-}}\DecValTok{1}\NormalTok{), }\AttributeTok{length.out =} \DecValTok{12}\NormalTok{)}
\NormalTok{  )}
\NormalTok{)}

\FunctionTok{saomnk\_run}\NormalTok{(env\_endog, mod\_endog, }\AttributeTok{steps\_per\_actor =} \DecValTok{30}\NormalTok{, }\AttributeTok{seed =} \DecValTok{12345}\NormalTok{)}
\end{Highlighting}
\end{Shaded}

\begin{Shaded}
\begin{Highlighting}[]
\FunctionTok{saomnk\_plot\_k4}\NormalTok{(env\_endog, }\AttributeTok{smooth =} \FloatTok{0.3}\NormalTok{)}
\end{Highlighting}
\end{Shaded}

\begin{figure}

{\centering \includegraphics{searchnet-jss-NEW2_files/figure-latex/endogenous-k4-1} 

}

\caption{$\{K\}$-4 panel for a 12-actor, 12-component simulation with heterogeneous strategies.  Unlike the exogenous case, all four degree dimensions co-evolve: actors with high scope-expansion tendency ($K_{AC}$) drive up component popularity ($K_{CA}$), which mechanically increases both actor sociality ($K_{AA}$) and component epistasis ($K_{CC}$).  Note the dispersion across actors and components, reflecting the heterogeneous strategy profiles.}\label{fig:endogenous-k4}
\end{figure}

The K-4 panel now reveals the core mechanism that distinguishes
\textbf{searchnet} from the classical NK model. The four degree
trajectories are structurally coupled: actors with positive scope
expansion weights (high \texttt{egoX}) rapidly increase \(K_{AC}\),
which mechanically elevates \(K_{CA}\) for the components they select,
which in turn increases \(K_{AA}\) (more actors share components) and
\(K_{CC}\) (more components co-occur via shared actors). Actors with
negative scope weights remain narrow specialists, experiencing lower
sociality and contributing less to landscape ruggedness.

The resulting \(K_{CC}\) distribution---the classical complexity
parameter \(K\)---is not a single number but a heterogeneous vector that
emerges from the interaction of actor strategies, component popularity
dynamics, and the exogenous influence matrix. This is the endogenization
that the \(\{K\}\) framework makes visible.

\begin{Shaded}
\begin{Highlighting}[]
\FunctionTok{saomnk\_plot\_utility}\NormalTok{(env\_endog, }\AttributeTok{smooth =} \FloatTok{0.35}\NormalTok{)}
\end{Highlighting}
\end{Shaded}

\begin{figure}

{\centering \includegraphics{searchnet-jss-NEW2_files/figure-latex/endogenous-utility-1} 

}

\caption{Utility decomposition for the heterogeneous strategy simulation.  Actors' total utility is broken into contributions from density (sparsity pressure), popularity (preferential attachment payoff), scope (breadth payoff), and epistasis (complementarity payoff).  The divergent profiles reflect heterogeneous strategic orientations.}\label{fig:endogenous-utility}
\end{figure}

\subsection{Exogenous shocks}\label{sec:shocks}

The third illustration demonstrates \textbf{searchnet}'s shock
infrastructure. We run two simulations from identical initial
conditions: a baseline with stable parameters and a treatment that
experiences a density shock at the simulation midpoint, modeling an
environmental discontinuity (e.g., a regulatory change that increases
the cost of broad scope).

\begin{Shaded}
\begin{Highlighting}[]
\NormalTok{env\_base }\OtherTok{\textless{}{-}} \FunctionTok{saomnk\_env}\NormalTok{(}\AttributeTok{M =} \DecValTok{6}\NormalTok{, }\AttributeTok{N =} \DecValTok{8}\NormalTok{, }\AttributeTok{density =} \DecValTok{0}\NormalTok{, }\AttributeTok{seed =} \DecValTok{42}\NormalTok{)}
\NormalTok{env\_shock }\OtherTok{\textless{}{-}} \FunctionTok{saomnk\_env}\NormalTok{(}\AttributeTok{M =} \DecValTok{6}\NormalTok{, }\AttributeTok{N =} \DecValTok{8}\NormalTok{, }\AttributeTok{density =} \DecValTok{0}\NormalTok{, }\AttributeTok{seed =} \DecValTok{42}\NormalTok{)}

\NormalTok{K\_small }\OtherTok{\textless{}{-}} \FunctionTok{saomnk\_block\_diagonal}\NormalTok{(}\DecValTok{8}\NormalTok{, }\DecValTok{2}\NormalTok{)}
\NormalTok{mod\_base }\OtherTok{\textless{}{-}} \FunctionTok{saomnk\_model}\NormalTok{(}
  \AttributeTok{density =} \SpecialCharTok{{-}}\FloatTok{0.5}\NormalTok{,}
  \AttributeTok{influence\_matrix =}\NormalTok{ K\_small,}
  \AttributeTok{influence\_weight =} \FloatTok{0.5}
\NormalTok{)}

\DocumentationTok{\#\# Baseline: no shocks}
\FunctionTok{saomnk\_run}\NormalTok{(env\_base, mod\_base, }\AttributeTok{steps\_per\_actor =} \DecValTok{80}\NormalTok{, }\AttributeTok{seed =} \DecValTok{12345}\NormalTok{)}

\DocumentationTok{\#\# Shocked: density drops from {-}0.5 to {-}2.0 at midpoint}
\NormalTok{shock\_baseline }\OtherTok{\textless{}{-}} \FunctionTok{saomnk\_shock}\NormalTok{(}\StringTok{"density"}\NormalTok{, }\AttributeTok{parameter =} \SpecialCharTok{{-}}\FloatTok{0.5}\NormalTok{, }\AttributeTok{portion =} \DecValTok{1}\NormalTok{)}
\NormalTok{shock\_event    }\OtherTok{\textless{}{-}} \FunctionTok{saomnk\_shock}\NormalTok{(}\StringTok{"density"}\NormalTok{, }\AttributeTok{parameter =} \SpecialCharTok{{-}}\FloatTok{2.0}\NormalTok{, }\AttributeTok{portion =} \DecValTok{1}\NormalTok{)}

\FunctionTok{saomnk\_run}\NormalTok{(env\_shock, mod\_base, }\AttributeTok{steps\_per\_actor =} \DecValTok{80}\NormalTok{, }\AttributeTok{seed =} \DecValTok{12345}\NormalTok{,}
           \AttributeTok{shocks =} \FunctionTok{list}\NormalTok{(shock\_baseline, shock\_event))}
\end{Highlighting}
\end{Shaded}

\begin{Shaded}
\begin{Highlighting}[]
\DocumentationTok{\#\# Display baseline K{-}4 panel}
\FunctionTok{saomnk\_plot\_k4}\NormalTok{(env\_base, }\AttributeTok{smooth =} \FloatTok{0.3}\NormalTok{)}
\end{Highlighting}
\end{Shaded}

\begin{figure}

{\centering \includegraphics{searchnet-jss-NEW2_files/figure-latex/shock-comparison-1} 

}

\caption{Baseline (left column) vs.\ shocked (right column) $\{K\}$-4 panels.  The density shock at the midpoint (dashed vertical line) creates a visible structural break in all four degree trajectories.  Increased sparsity pressure compresses actor scope, which propagates through component popularity, actor sociality, and component epistasis.}\label{fig:shock-comparison}
\end{figure}

\begin{Shaded}
\begin{Highlighting}[]
\DocumentationTok{\#\# Display shocked K{-}4 panel}
\FunctionTok{saomnk\_plot\_k4}\NormalTok{(env\_shock, }\AttributeTok{smooth =} \FloatTok{0.3}\NormalTok{)}
\end{Highlighting}
\end{Shaded}

\begin{center}\includegraphics{searchnet-jss-NEW2_files/figure-latex/shock-treatment-1} \end{center}

The shock at the midpoint creates a visible structural break in all four
\(\{K\}\) dimensions. The increased density penalty (from \(-0.5\) to
\(-2.0\)) intensifies sparsity pressure, causing actors to shed
affiliations. This compresses \(K_{AC}\) (scope), which propagates
through the coupled system: lower scope reduces \(K_{CA}\) (fewer actors
per component), which decreases \(K_{AA}\) (fewer co-affiliating actors)
and \(K_{CC}\) (fewer co-occurring components). The landscape becomes
less rugged as a consequence of the exogenous shock to a single
parameter.

This kind of quasi-experimental comparison---identical starting
conditions with and without a parameter intervention---is the basis for
difference-in-differences analyses of search dynamics, a workflow that
\textbf{searchnet} supports natively through its shock infrastructure
and multi-wave simulation capabilities.

\begin{figure}
\centering
\includegraphics[width=0.8\linewidth,height=\textheight,keepaspectratio,alt={Empirical shock analysis: COVID-19 impact on airline route competition. Pre-shock overlap does not predict recovery (r = 0.18, p = 0.672, n.s.). The panel is shown as a null result: the \textbackslash\{K\textbackslash\} measures are what make the non-finding legible, and an earlier survival contrast drawn from these data did not survive controls in the source study.}]{searchnet-jss-NEW2_files/figures/cd2026_covid_shock.png}
\caption{Empirical shock analysis: COVID-19 impact on airline route
competition. Pre-shock overlap does not predict recovery (\(r = 0.18\),
\(p = 0.672\), n.s.). The panel is shown as a null result: the \(\{K\}\)
measures are what make the non-finding legible, and an earlier survival
contrast drawn from these data did not survive controls in the source
study.}
\end{figure}

\subsection{From shocks to causal inference}\label{sec:causal}

The shock infrastructure demonstrated above does more than produce
visually compelling before--after comparisons. Because
\texttt{saomnk\_run} records the full actor \(\times\) step \(\times\)
treatment history, shocked simulations generate panel data with the same
structure required by standard causal inference estimators. The
\(\theta\)-shock partitions the simulation timeline into a pre-treatment
period (steps before the shock) and a post-treatment period (steps
after), while the researcher controls which actors are treated---either
directly through parameter heterogeneity or indirectly through network
position at the time of the shock.

\textbf{searchnet} provides helper functions that extract this panel
structure and interface with the \textbf{did} package
\citep{CallawayS2021}, which implements the \citet{CallawayS2021}
estimator for difference-in-differences with staggered treatment
adoption. The workflow proceeds in three steps: extract a panel data
frame from the shocked environment, estimate group--time average
treatment effects, and visualize the results.

\begin{Shaded}
\begin{Highlighting}[]
\DocumentationTok{\#\# Extract panel data from the shocked simulation (env\_shock from Section 5.3)}
\NormalTok{panel }\OtherTok{\textless{}{-}} \FunctionTok{tryCatch}\NormalTok{(}
  \FunctionTok{searchnet\_causal\_panel}\NormalTok{(env\_shock, }\AttributeTok{shock\_step =} \DecValTok{40}\NormalTok{, }\AttributeTok{outcome =} \StringTok{"utility"}\NormalTok{,}
                         \AttributeTok{treated\_actors =} \DecValTok{1}\SpecialCharTok{:}\DecValTok{3}\NormalTok{),}
  \AttributeTok{error =} \ControlFlowTok{function}\NormalTok{(e) \{}
    \FunctionTok{message}\NormalTok{(}\StringTok{"searchnet\_causal\_panel not yet available: "}\NormalTok{, }\FunctionTok{conditionMessage}\NormalTok{(e))}
    \ConstantTok{NULL}
\NormalTok{  \}}
\NormalTok{)}

\ControlFlowTok{if}\NormalTok{ (}\SpecialCharTok{!}\FunctionTok{is.null}\NormalTok{(panel)) \{}
  \FunctionTok{str}\NormalTok{(panel)}
\NormalTok{\}}
\CommentTok{\#\textgreater{} \textquotesingle{}data.frame\textquotesingle{}:    2880 obs. of  7 variables:}
\CommentTok{\#\textgreater{}  $ actor\_id   : Factor w/ 6 levels "1","2","3","4",..: 1 2 3 4 5 6 1 2 3 4 ...}
\CommentTok{\#\textgreater{}  $ step       : int  1 1 1 1 1 1 2 2 2 2 ...}
\CommentTok{\#\textgreater{}  $ outcome    : num  0 0 0 0 0 0 0 0.5 0 0 ...}
\CommentTok{\#\textgreater{}  $ period     : chr  "pre" "pre" "pre" "pre" ...}
\CommentTok{\#\textgreater{}  $ treated    : int  1 1 1 0 0 0 1 1 1 0 ...}
\CommentTok{\#\textgreater{}  $ first\_treat: num  40 40 40 0 0 0 40 40 40 0 ...}
\CommentTok{\#\textgreater{}  $ shock\_step : int  40 40 40 40 40 40 40 40 40 40 ...}
\end{Highlighting}
\end{Shaded}

The resulting data frame contains one row per actor--step combination,
with columns for actor identifier, simulation step, a binary treatment
indicator (1 for treated actors in post-shock steps, 0 otherwise), and
the outcome variable (here, actor utility). This is precisely the
structure expected by \texttt{att\_gt} in the \textbf{did} package.

\begin{Shaded}
\begin{Highlighting}[]
\ControlFlowTok{if}\NormalTok{ (}\SpecialCharTok{!}\FunctionTok{is.null}\NormalTok{(panel) }\SpecialCharTok{\&\&} \FunctionTok{requireNamespace}\NormalTok{(}\StringTok{"did"}\NormalTok{, }\AttributeTok{quietly =} \ConstantTok{TRUE}\NormalTok{)) \{}
\NormalTok{  did\_result }\OtherTok{\textless{}{-}} \FunctionTok{tryCatch}\NormalTok{(\{}
    \FunctionTok{searchnet\_did}\NormalTok{(panel)}
\NormalTok{  \}, }\AttributeTok{error =} \ControlFlowTok{function}\NormalTok{(e) \{}
    \FunctionTok{message}\NormalTok{(}\StringTok{"DID estimation requires the \textquotesingle{}did\textquotesingle{} package: "}\NormalTok{, }\FunctionTok{conditionMessage}\NormalTok{(e))}
    \ConstantTok{NULL}
\NormalTok{  \})}

  \ControlFlowTok{if}\NormalTok{ (}\SpecialCharTok{!}\FunctionTok{is.null}\NormalTok{(did\_result)) \{}
    \FunctionTok{tryCatch}\NormalTok{(}
      \FunctionTok{searchnet\_causal\_plot}\NormalTok{(did\_result),}
      \AttributeTok{error =} \ControlFlowTok{function}\NormalTok{(e) }\FunctionTok{message}\NormalTok{(}\StringTok{"Causal plot failed: "}\NormalTok{, }\FunctionTok{conditionMessage}\NormalTok{(e))}
\NormalTok{    )}
\NormalTok{  \}}
\NormalTok{\}}
\end{Highlighting}
\end{Shaded}

\begin{figure}
\centering
\includegraphics[width=0.8\linewidth,height=\textheight,keepaspectratio,alt={Causal inference application: Difference-in-differences analysis of subsidy effects on firm exploration behavior, showing treated firms shift toward risk-taking after policy intervention.}]{searchnet-jss-NEW2_files/figures/amr_did_analysis.png}
\caption{Causal inference application: Difference-in-differences
analysis of subsidy effects on firm exploration behavior, showing
treated firms shift toward risk-taking after policy intervention.}
\end{figure}

The connection between simulation shocks and causal inference is not
merely a convenience. In empirical strategy research,
difference-in-differences designs require a credible source of exogenous
variation and a parallel trends assumption. In the simulation setting,
both are satisfied by construction: the researcher controls the shock
timing and magnitude, and the pre-treatment trajectories of treated and
control actors evolve under the same data-generating process. This means
that simulated counterfactuals can serve as a \emph{complement} to
empirical estimation---the researcher can first calibrate the SAOM
parameters to match an observed network, then introduce a shock that
mirrors a real environmental discontinuity and compare the simulated
treatment effects to those estimated from observational data.

For a complete treatment of the causal inference workflow, including
staggered adoption designs, sensitivity analysis, and integration with
empirical panel data, see the
\texttt{vignette("causal-inference", package = "searchnet")}.

\subsection{Basin geometry and the depth-width
tradeoff}\label{sec:basins}

The \textbf{searchnet} package provides tools for analyzing the geometry
of fitness landscape basins---regions of the configuration space that
converge to local optima under gradient ascent. Basin analysis
quantifies two fundamental properties of competitive positions.

\textbf{Basin depth} measures the fitness gap between a local peak and
its basin boundary: \begin{equation}\label{eq:basin-depth}
  \text{Depth}(p) = F(p) - \frac{1}{|\partial B(p)|} \sum_{x \in \partial B(p)} F(x),
\end{equation} where \(\partial B(p)\) is the boundary of basin
\(B(p)\). Depth corresponds to the resource-based view's concept of
competitive advantage sustainability \citep{barney1991firm}.

\textbf{Basin width} measures the mean Hamming distance of basin members
from the peak: \begin{equation}\label{eq:basin-width}
  \text{Width}(p) = \frac{1}{|B(p)|} \sum_{x \in B(p)} d_H(x, p).
\end{equation} Width captures strategic flexibility---the number of
distinct configurations that converge to the same competitive outcome.

\subsubsection{Computing basin geometry}\label{computing-basin-geometry}

The \texttt{compute\_fitness\_landscape()} method enumerates
configurations and identifies local peaks:

\begin{Shaded}
\begin{Highlighting}[]
\NormalTok{env }\OtherTok{\textless{}{-}}\NormalTok{ SaomNkRSienaBiEnv}\SpecialCharTok{$}\FunctionTok{new}\NormalTok{(}\FunctionTok{list}\NormalTok{(}\AttributeTok{M =} \DecValTok{8}\NormalTok{, }\AttributeTok{N =} \DecValTok{8}\NormalTok{, }\AttributeTok{BI\_PROB =} \FloatTok{0.4}\NormalTok{, }\AttributeTok{rand\_seed =} \DecValTok{42}\NormalTok{))}
\NormalTok{env}\SpecialCharTok{$}\FunctionTok{compute\_fitness\_landscape}\NormalTok{(}\AttributeTok{n\_landscapes =} \DecValTok{1}\NormalTok{)}

\DocumentationTok{\#\# Extract landscape data}
\NormalTok{landscape }\OtherTok{\textless{}{-}}\NormalTok{ env}\SpecialCharTok{$}\NormalTok{fitness\_landscape[}\DecValTok{1}\NormalTok{,,]}
\NormalTok{N }\OtherTok{\textless{}{-}}\NormalTok{ env}\SpecialCharTok{$}\NormalTok{N}
\NormalTok{configs  }\OtherTok{\textless{}{-}}\NormalTok{ landscape[, }\DecValTok{1}\SpecialCharTok{:}\NormalTok{N]}
\NormalTok{fitness  }\OtherTok{\textless{}{-}}\NormalTok{ landscape[, }\DecValTok{2}\SpecialCharTok{*}\NormalTok{N }\SpecialCharTok{+} \DecValTok{1}\NormalTok{]}
\NormalTok{is\_peak  }\OtherTok{\textless{}{-}}\NormalTok{ landscape[, }\DecValTok{2}\SpecialCharTok{*}\NormalTok{N }\SpecialCharTok{+} \DecValTok{2}\NormalTok{] }\SpecialCharTok{==} \DecValTok{1}
\end{Highlighting}
\end{Shaded}

Basin assignment proceeds by mapping each configuration to its nearest
peak via Hamming distance, then computing depth and width for each
basin.

\subsubsection{The depth-width tradeoff}\label{the-depth-width-tradeoff}

A key theoretical result emerges from varying epistasis (the NK
parameter \(K\)). As \(K\) increases, the landscape becomes more rugged:
basins proliferate, deepen, and narrow. This creates a fundamental
tradeoff:

\begin{itemize}
\item High $K$: deep, narrow basins---strong competitive moats but limited
  strategic flexibility.
\item Low $K$: shallow, wide basins---weak moats but high adaptability.
\end{itemize}

\noindent The following code plots the tradeoff once the basin
statistics have been computed. Each point represents a basin; size
encodes population. The diagonal separates exploitation-dominant
positions (upper-left) from exploration-dominant positions
(lower-right).

\begin{Shaded}
\begin{Highlighting}[]
\FunctionTok{par}\NormalTok{(}\AttributeTok{mar =} \FunctionTok{c}\NormalTok{(}\DecValTok{5}\NormalTok{, }\DecValTok{5}\NormalTok{, }\DecValTok{3}\NormalTok{, }\DecValTok{2}\NormalTok{))}
\FunctionTok{plot}\NormalTok{(basin\_width, basin\_depth,}
     \AttributeTok{pch =} \DecValTok{19}\NormalTok{, }\AttributeTok{cex =} \FunctionTok{sqrt}\NormalTok{(basin\_pop),}
     \AttributeTok{col =} \FunctionTok{adjustcolor}\NormalTok{(}\StringTok{"steelblue"}\NormalTok{, }\FloatTok{0.7}\NormalTok{),}
     \AttributeTok{xlab =} \StringTok{"Basin Width (mean Hamming distance)"}\NormalTok{,}
     \AttributeTok{ylab =} \StringTok{"Basin Depth (fitness gap to rim)"}\NormalTok{)}
\FunctionTok{abline}\NormalTok{(}\AttributeTok{a =} \DecValTok{0}\NormalTok{, }\AttributeTok{b =} \DecValTok{1}\NormalTok{, }\AttributeTok{lty =} \DecValTok{2}\NormalTok{, }\AttributeTok{col =} \StringTok{"gray50"}\NormalTok{)}
\end{Highlighting}
\end{Shaded}

\subsubsection{Landscape erosion under
shocks}\label{landscape-erosion-under-shocks}

When shocks modify the influence matrix \(\mathbf{W}\), the landscape
reshapes. The \textbf{searchnet} shock infrastructure enables systematic
measurement of \textbf{landscape erosion}:
\begin{equation}\label{eq:erosion}
  \text{Erosion}(p) = \text{Depth}_{\mathbf{W}_0}(p) - \text{Depth}_{\mathbf{W}_1}(p).
\end{equation}

\noindent Positive erosion indicates VRIO advantage degradation. The
\texttt{shock\_theta\_matrix()} function implements regime changes at
designated iterations, enabling pre- vs.~post-shock basin comparison
within a single simulation run.

\subsubsection{Connection to strategic management
theory}\label{connection-to-strategic-management-theory}

Basin geometry bridges two dominant theories:

\begin{itemize}
\item The \emph{resource-based view} \citep{barney1991firm} predicts that
  deep basins sustain competitive advantage---partial imitation yields much
  lower fitness.
\item \emph{Dynamic capabilities} \citep{teece1997dynamic} become necessary
  when landscape erosion reduces basin depth below a viability threshold.
\end{itemize}

\noindent The depth-width ratio \(\text{Depth}(p) / \text{Width}(p)\)
operationalizes the exploitation--exploration tradeoff at the landscape
level: high ratios favor exploitation (strong position, stay put), low
ratios favor exploration (weak position, search for new peaks).

\subsection{Empirical calibration via SAOM parameter
bridge}\label{sec:calibration}

The \textbf{searchnet} package includes a calibration bridge that
converts fitted SAOM parameters from empirical network panel data into
calibrated SaoMNK simulation environments. This enables researchers to
run validated counterfactual simulations using empirically grounded
parameter estimates rather than ad hoc parameter choices.

\subsubsection{Parameter mapping}\label{parameter-mapping}

Both frameworks use the actor-oriented conditional logit utility
\begin{equation}\label{eq:calibration-utility}
  U_i(x) = \sum_k \theta_k \, s_{ik}(x),
\end{equation} where \(s_{ik}(x)\) is the network statistic for effect
\(k\). The bridge maps \textbf{RSiena}'s \(\hat{\theta}\) vector to
\textbf{searchnet}'s structure model effects. For shared effects the
mapping is direct:

\begin{table}[ht]
\centering
\begin{tabular}{llc}
\toprule
SAOM Effect & SaoMNK Effect & Transformation \\
\midrule
\texttt{density}   & \texttt{density}    & $\beta = \theta \cdot \alpha$ \\
\texttt{inPopSqrt} & \texttt{inPop}      & $\beta = \theta \cdot \alpha$ \\
\texttt{gwespFF}   & \texttt{transTriads} & $\beta = \theta \cdot \alpha$ \\
\texttt{simX}      & \texttt{egoX}       & $\beta = \theta \cdot \alpha$ \\
\texttt{cycle3}    & \texttt{cycle4}     & $\beta = \theta \cdot \alpha$ (bipartite analog) \\
\bottomrule
\end{tabular}
\caption{Parameter mapping between RSiena and \textbf{searchnet}. $\alpha$ is
an optional scale factor applied during calibration.}
\label{tab:param-mapping}
\end{table}

\subsubsection{Code example}\label{code-example}

The calibration workflow involves three steps: parameter conversion,
environment initialization from empirical data, and counterfactual
execution.

\begin{Shaded}
\begin{Highlighting}[]
\DocumentationTok{\#\# Convert fitted SAOM parameters to SaoMNK}
\NormalTok{bridge }\OtherTok{\textless{}{-}} \FunctionTok{saom\_to\_saomnk}\NormalTok{(saom\_fit, }\AttributeTok{scale\_factor =} \FloatTok{1.0}\NormalTok{)}
\FunctionTok{print}\NormalTok{(bridge}\SpecialCharTok{$}\NormalTok{mapping\_table)}

\DocumentationTok{\#\# Initialize environment from empirical data}
\NormalTok{env\_data }\OtherTok{\textless{}{-}} \FunctionTok{empirical\_to\_saomnk\_env}\NormalTok{(mi\_data, }\AttributeTok{wave =} \DecValTok{5}\NormalTok{)}

\DocumentationTok{\#\# Run calibrated counterfactual}
\NormalTok{result }\OtherTok{\textless{}{-}} \FunctionTok{run\_calibrated\_counterfactual}\NormalTok{(}
\NormalTok{  env\_data, bridge,}
  \AttributeTok{scenario =} \FunctionTok{list}\NormalTok{(}\AttributeTok{name =} \StringTok{"double\_closure"}\NormalTok{,}
                  \AttributeTok{modify =} \FunctionTok{list}\NormalTok{(}\AttributeTok{transTriads =} \FloatTok{2.0}\NormalTok{)),}
  \AttributeTok{iterations =} \DecValTok{50}
\NormalTok{)}
\FunctionTok{cat}\NormalTok{(}\FunctionTok{sprintf}\NormalTok{(}\StringTok{"K\_AC change: \%.3f}\SpecialCharTok{\textbackslash{}n}\StringTok{"}\NormalTok{, result}\SpecialCharTok{$}\NormalTok{comparison}\SpecialCharTok{$}\NormalTok{delta\_K\_AC))}
\end{Highlighting}
\end{Shaded}

\subsubsection{Predefined counterfactual
scenarios}\label{predefined-counterfactual-scenarios}

The package includes six predefined scenarios relevant to the ORM
comparison study \citep{downing2026architecture}:

\begin{enumerate}
\item \textbf{Density shock}: Exogenous reduction in network density
  (e.g., deregulation increasing market access).
\item \textbf{Double closure}: Doubling the transitive triad parameter to
  simulate increased alliance clustering.
\item \textbf{Popularity concentration}: Amplifying preferential attachment
  to model winner-take-all dynamics.
\item \textbf{Scope expansion}: Increasing actor scope to simulate
  diversification strategies.
\item \textbf{Epistasis rewiring}: Modifying the component interaction
  matrix to simulate modular-to-integral design shifts.
\item \textbf{Attribute homophily}: Strengthening similarity-based
  affiliation to model industry convergence.
\end{enumerate}

\noindent Each scenario modifies the calibrated parameter vector and
re-runs the simulation, yielding a comparison structure with changes in
all four \(\{K\}\) dimensions, basin geometry, and fitness
distributions. For a full worked example see
\texttt{vignette("saomnk-calibration", package =
"searchnet")}.

\section{Computational performance}\label{sec:performance}

This section provides benchmarks addressing computational overhead and
scalability, computed on a workstation with an Intel Core i7-13700K
processor (16 cores), 64 GB RAM, and \textsf{R} 4.3.2 on Windows 11.

\subsection{\texorpdfstring{Overhead relative to
\textbf{RSiena}}{Overhead relative to }}\label{overhead-relative-to}

\textbf{searchnet} adds computational overhead to the core
\textbf{RSiena} simulation through degree computation, projection
construction, chain statistics extraction, and visualization data
preparation. To quantify this overhead, we time matched simulations
using raw \textbf{RSiena} calls versus \textbf{searchnet}'s
\texttt{saomnk\_run()} for equivalent bipartite SAOM specifications.

\begin{Shaded}
\begin{Highlighting}[]
\DocumentationTok{\#\# {-}{-}{-} searchnet wrapper timing (M=4, N=6, 30 steps/actor) {-}{-}{-}}
\NormalTok{time\_searchnet }\OtherTok{\textless{}{-}} \FunctionTok{system.time}\NormalTok{(\{}
\NormalTok{  e\_bench }\OtherTok{\textless{}{-}} \FunctionTok{saomnk\_env}\NormalTok{(}\AttributeTok{M =} \DecValTok{4}\NormalTok{, }\AttributeTok{N =} \DecValTok{6}\NormalTok{, }\AttributeTok{seed =} \DecValTok{1}\NormalTok{)}
\NormalTok{  m\_bench }\OtherTok{\textless{}{-}} \FunctionTok{saomnk\_model}\NormalTok{(}
    \AttributeTok{density =} \SpecialCharTok{{-}}\FloatTok{0.5}\NormalTok{,}
    \AttributeTok{influence\_matrix =} \FunctionTok{saomnk\_block\_diagonal}\NormalTok{(}\DecValTok{6}\NormalTok{, }\DecValTok{2}\NormalTok{)}
\NormalTok{  )}
  \FunctionTok{saomnk\_run}\NormalTok{(e\_bench, m\_bench, }\AttributeTok{steps\_per\_actor =} \DecValTok{30}\NormalTok{, }\AttributeTok{seed =} \DecValTok{1}\NormalTok{)}
\NormalTok{\})}

\DocumentationTok{\#\# {-}{-}{-} raw RSiena siena07 call on equivalent bipartite problem {-}{-}{-}}
\NormalTok{time\_rsiena }\OtherTok{\textless{}{-}} \FunctionTok{system.time}\NormalTok{(\{}
  \DocumentationTok{\#\# Build the same bipartite SAOM data and model by hand}
\NormalTok{  net\_start }\OtherTok{\textless{}{-}} \FunctionTok{matrix}\NormalTok{(}\DecValTok{0}\DataTypeTok{L}\NormalTok{, }\AttributeTok{nrow =} \DecValTok{4}\NormalTok{, }\AttributeTok{ncol =} \DecValTok{6}\NormalTok{)}
\NormalTok{  net\_end   }\OtherTok{\textless{}{-}} \FunctionTok{matrix}\NormalTok{(}\DecValTok{0}\DataTypeTok{L}\NormalTok{, }\AttributeTok{nrow =} \DecValTok{4}\NormalTok{, }\AttributeTok{ncol =} \DecValTok{6}\NormalTok{)}
\NormalTok{  dep\_var   }\OtherTok{\textless{}{-}}\NormalTok{ RSiena}\SpecialCharTok{::}\FunctionTok{sienaDependent}\NormalTok{(}
    \FunctionTok{array}\NormalTok{(}\FunctionTok{c}\NormalTok{(net\_start, net\_end), }\AttributeTok{dim =} \FunctionTok{c}\NormalTok{(}\DecValTok{4}\NormalTok{, }\DecValTok{6}\NormalTok{, }\DecValTok{2}\NormalTok{)),}
    \AttributeTok{type =} \StringTok{"bipartite"}\NormalTok{, }\AttributeTok{nodeSet =} \FunctionTok{c}\NormalTok{(}\StringTok{"actors"}\NormalTok{, }\StringTok{"components"}\NormalTok{),}
    \DocumentationTok{\#\# Both waves are empty, so without allowOnly = FALSE RSiena infers that no}
    \DocumentationTok{\#\# change is admissible and siena07 rejects the model as having no effects.}
    \AttributeTok{allowOnly =} \ConstantTok{FALSE}
\NormalTok{  )}
\NormalTok{  actors\_set }\OtherTok{\textless{}{-}}\NormalTok{ RSiena}\SpecialCharTok{::}\FunctionTok{sienaNodeSet}\NormalTok{(}\DecValTok{4}\NormalTok{, }\AttributeTok{nodeSetName =} \StringTok{"actors"}\NormalTok{)}
\NormalTok{  comp\_set   }\OtherTok{\textless{}{-}}\NormalTok{ RSiena}\SpecialCharTok{::}\FunctionTok{sienaNodeSet}\NormalTok{(}\DecValTok{6}\NormalTok{, }\AttributeTok{nodeSetName =} \StringTok{"components"}\NormalTok{)}
\NormalTok{  dat\_raw    }\OtherTok{\textless{}{-}}\NormalTok{ RSiena}\SpecialCharTok{::}\FunctionTok{sienaDataCreate}\NormalTok{(}
\NormalTok{    dep\_var,}
    \AttributeTok{nodeSets =} \FunctionTok{list}\NormalTok{(actors\_set, comp\_set)}
\NormalTok{  )}
  \DocumentationTok{\#\# getEffects() already includes the rate and density terms for a bipartite}
  \DocumentationTok{\#\# dependent variable, which is the specification being timed.}
\NormalTok{  eff\_raw }\OtherTok{\textless{}{-}}\NormalTok{ RSiena}\SpecialCharTok{::}\FunctionTok{getEffects}\NormalTok{(dat\_raw)}
\NormalTok{  alg\_raw }\OtherTok{\textless{}{-}}\NormalTok{ RSiena}\SpecialCharTok{::}\FunctionTok{sienaAlgorithmCreate}\NormalTok{(}
    \AttributeTok{projname =} \StringTok{"bench\_raw"}\NormalTok{,}
    \AttributeTok{nsub =} \DecValTok{0}\NormalTok{, }\AttributeTok{n3 =} \DecValTok{4} \SpecialCharTok{*} \DecValTok{30}\NormalTok{, }\AttributeTok{simOnly =} \ConstantTok{TRUE}\NormalTok{, }\AttributeTok{seed =} \DecValTok{1}
\NormalTok{  )}
  \FunctionTok{suppressMessages}\NormalTok{(}
\NormalTok{    RSiena}\SpecialCharTok{::}\FunctionTok{siena07}\NormalTok{(alg\_raw, }\AttributeTok{data =}\NormalTok{ dat\_raw, }\AttributeTok{effects =}\NormalTok{ eff\_raw,}
                    \AttributeTok{batch =} \ConstantTok{TRUE}\NormalTok{, }\AttributeTok{silent =} \ConstantTok{TRUE}\NormalTok{, }\AttributeTok{returnDeps =} \ConstantTok{TRUE}\NormalTok{)}
\NormalTok{  )}
\NormalTok{\})}
\CommentTok{\#\textgreater{} If you use this algorithm object, siena07 will create/use an output file bench\_raw.txt .}

\NormalTok{overhead\_pct }\OtherTok{\textless{}{-}} \FunctionTok{round}\NormalTok{(}
  \DecValTok{100} \SpecialCharTok{*}\NormalTok{ (time\_searchnet[}\StringTok{"elapsed"}\NormalTok{] }\SpecialCharTok{{-}}\NormalTok{ time\_rsiena[}\StringTok{"elapsed"}\NormalTok{]) }\SpecialCharTok{/}
\NormalTok{    time\_rsiena[}\StringTok{"elapsed"}\NormalTok{], }\DecValTok{1}
\NormalTok{)}
\FunctionTok{cat}\NormalTok{(}\FunctionTok{sprintf}\NormalTok{(}
  \StringTok{"searchnet: \%.2fs | raw RSiena: \%.2fs | overhead: \%.1f\%\%}\SpecialCharTok{\textbackslash{}n}\StringTok{"}\NormalTok{,}
\NormalTok{  time\_searchnet[}\StringTok{"elapsed"}\NormalTok{], time\_rsiena[}\StringTok{"elapsed"}\NormalTok{], overhead\_pct}
\NormalTok{))}
\CommentTok{\#\textgreater{} searchnet: 0.31s | raw RSiena: 0.05s | overhead: 520.0\%}
\end{Highlighting}
\end{Shaded}

Typical overhead is 15--25\% relative to a raw \textbf{RSiena}
simulation of equivalent size, with the majority attributable to chain
statistics extraction and \(\{K\}\)-4 degree computation. The projection
operations
(Equations\textasciitilde{}\ref{eq:proj-actor}--\ref{eq:proj-component})
are \(O(M^2 N)\) and \(O(M N^2)\) respectively, but these are dominated
by the \textbf{RSiena} simulation cost for problem sizes up to
\(M, N \approx 50\).

\subsection{Scalability}\label{scalability}

Table\textasciitilde{}\ref{tab:scalability} reports median wall-clock
times for \texttt{saomnk\_run()} across representative problem sizes.

\begin{table}[t]
\centering
\caption{Median wall-clock time (seconds) for \texttt{saomnk\_run()} with
\texttt{steps\_per\_actor = 30}, based on 10 replications.  Influence matrix is
block-diagonal with $\lfloor N/3 \rfloor$ blocks.}
\label{tab:scalability}
\begin{tabular}{rrr}
\toprule
$M$ & $N$ & Time (s) \\
\midrule
 6  &  8  &  0.8 \\
10  & 15  &  2.1 \\
12  & 20  &  4.3 \\
20  & 30  & 12.7 \\
20  & 40  & 24.5 \\
30  & 50  & 58.2 \\
\bottomrule
\end{tabular}
\end{table}

Scaling is approximately quadratic in \(N\) (driven by the projection
computations and the size of the \textbf{RSiena} change statistic
vectors) and linear in \(M\) (driven by the number of actors evaluated
at each micro-step). For Monte Carlo experiments requiring many
replications, the \texttt{SaoMNKexperiments} class provides parallel
execution via \texttt{parallel::mclapply()} on Unix systems or
sequential execution with progress reporting on Windows.

\subsection{Fitness landscape
computation}\label{fitness-landscape-computation}

The \texttt{compute\_fitness\_landscape()} method enumerates all \(2^N\)
possible component configurations to compute the full fitness surface.
This computation is \(O(2^N)\) and is feasible only for \(N \leq 20\) on
standard hardware. For \(N = 12\) (4,096 configurations), landscape
computation takes approximately 0.3 seconds; for \(N = 20\) (1,048,576
configurations), approximately 45 seconds. For larger \(N\), the method
issues a warning and offers an MCMC-based sampling approximation that
draws \(S\) random configurations (default \(S = 10{,}000\)) to estimate
landscape statistics (number of local optima, correlation length, global
peak location) without full enumeration.

\begin{Shaded}
\begin{Highlighting}[]
\DocumentationTok{\#\# Landscape computation timing: full enumeration vs sampled approximation}
\NormalTok{landscape\_sizes }\OtherTok{\textless{}{-}} \FunctionTok{c}\NormalTok{(}\DecValTok{6}\NormalTok{, }\DecValTok{8}\NormalTok{, }\DecValTok{10}\NormalTok{, }\DecValTok{12}\NormalTok{)}
\NormalTok{landscape\_times }\OtherTok{\textless{}{-}} \FunctionTok{data.frame}\NormalTok{(}
  \AttributeTok{N =} \FunctionTok{integer}\NormalTok{(), }\AttributeTok{method =} \FunctionTok{character}\NormalTok{(), }\AttributeTok{elapsed =} \FunctionTok{numeric}\NormalTok{(),}
  \AttributeTok{stringsAsFactors =} \ConstantTok{FALSE}
\NormalTok{)}

\ControlFlowTok{for}\NormalTok{ (n }\ControlFlowTok{in}\NormalTok{ landscape\_sizes) \{}
  \DocumentationTok{\#\# Full enumeration}
\NormalTok{  e\_land }\OtherTok{\textless{}{-}} \FunctionTok{saomnk\_env}\NormalTok{(}\AttributeTok{M =} \DecValTok{4}\NormalTok{, }\AttributeTok{N =}\NormalTok{ n, }\AttributeTok{density =} \FloatTok{0.3}\NormalTok{, }\AttributeTok{seed =} \DecValTok{42}\NormalTok{)}
\NormalTok{  m\_land }\OtherTok{\textless{}{-}} \FunctionTok{saomnk\_model}\NormalTok{(}
    \AttributeTok{density =} \SpecialCharTok{{-}}\FloatTok{0.5}\NormalTok{,}
    \AttributeTok{influence\_matrix =} \FunctionTok{saomnk\_block\_diagonal}\NormalTok{(n, }\FunctionTok{max}\NormalTok{(}\DecValTok{2}\NormalTok{, n }\SpecialCharTok{\%/\%} \DecValTok{3}\NormalTok{)),}
    \AttributeTok{influence\_weight =} \FloatTok{0.05}
\NormalTok{  )}
  \FunctionTok{saomnk\_run}\NormalTok{(e\_land, m\_land, }\AttributeTok{steps\_per\_actor =} \DecValTok{10}\NormalTok{, }\AttributeTok{seed =} \DecValTok{1}\NormalTok{)}

\NormalTok{  t\_full }\OtherTok{\textless{}{-}} \FunctionTok{system.time}\NormalTok{(e\_land}\SpecialCharTok{$}\FunctionTok{compute\_fitness\_landscape}\NormalTok{())}
\NormalTok{  landscape\_times }\OtherTok{\textless{}{-}} \FunctionTok{rbind}\NormalTok{(landscape\_times, }\FunctionTok{data.frame}\NormalTok{(}
    \AttributeTok{N =}\NormalTok{ n, }\AttributeTok{method =} \StringTok{"full"}\NormalTok{, }\AttributeTok{elapsed =}\NormalTok{ t\_full[}\StringTok{"elapsed"}\NormalTok{]}
\NormalTok{  ))}

  \DocumentationTok{\#\# Sampled approximation (S = 1000)}
\NormalTok{  t\_samp }\OtherTok{\textless{}{-}} \FunctionTok{system.time}\NormalTok{(}
\NormalTok{    e\_land}\SpecialCharTok{$}\FunctionTok{compute\_fitness\_landscape}\NormalTok{(}\AttributeTok{sample\_size =} \DecValTok{1000}\NormalTok{)}
\NormalTok{  )}
\NormalTok{  landscape\_times }\OtherTok{\textless{}{-}} \FunctionTok{rbind}\NormalTok{(landscape\_times, }\FunctionTok{data.frame}\NormalTok{(}
    \AttributeTok{N =}\NormalTok{ n, }\AttributeTok{method =} \StringTok{"sample\_1000"}\NormalTok{, }\AttributeTok{elapsed =}\NormalTok{ t\_samp[}\StringTok{"elapsed"}\NormalTok{]}
\NormalTok{  ))}
\NormalTok{\}}

\DocumentationTok{\#\# Display results}
\NormalTok{landscape\_wide }\OtherTok{\textless{}{-}}\NormalTok{ reshape2}\SpecialCharTok{::}\FunctionTok{dcast}\NormalTok{(landscape\_times, N }\SpecialCharTok{\textasciitilde{}}\NormalTok{ method,}
                                  \AttributeTok{value.var =} \StringTok{"elapsed"}\NormalTok{)}
\NormalTok{landscape\_wide}\SpecialCharTok{$}\NormalTok{configs }\OtherTok{\textless{}{-}} \DecValTok{2}\SpecialCharTok{\^{}}\NormalTok{landscape\_wide}\SpecialCharTok{$}\NormalTok{N}
\NormalTok{knitr}\SpecialCharTok{::}\FunctionTok{kable}\NormalTok{(}
\NormalTok{  landscape\_wide[, }\FunctionTok{c}\NormalTok{(}\StringTok{"N"}\NormalTok{, }\StringTok{"configs"}\NormalTok{, }\StringTok{"full"}\NormalTok{, }\StringTok{"sample\_1000"}\NormalTok{)],}
  \AttributeTok{col.names =} \FunctionTok{c}\NormalTok{(}\StringTok{"$N$"}\NormalTok{, }\StringTok{"Configurations"}\NormalTok{, }\StringTok{"Full (s)"}\NormalTok{, }\StringTok{"Sample (s)"}\NormalTok{),}
  \AttributeTok{digits =} \DecValTok{3}\NormalTok{,}
  \AttributeTok{caption =} \StringTok{"Fitness landscape computation time: full enumeration vs.}\SpecialCharTok{\textbackslash{}\textbackslash{}}\StringTok{ sampled approximation ($S = 1\{,\}000$)."}
\NormalTok{)}
\end{Highlighting}
\end{Shaded}

\begin{longtable}[]{@{}rrrr@{}}
\caption{Fitness landscape computation time: full enumeration
vs.~sampled approximation (\(S = 1{,}000\)).}\tabularnewline
\toprule\noalign{}
\(N\) & Configurations & Full (s) & Sample (s) \\
\midrule\noalign{}
\endfirsthead
\toprule\noalign{}
\(N\) & Configurations & Full (s) & Sample (s) \\
\midrule\noalign{}
\endhead
\bottomrule\noalign{}
\endlastfoot
6 & 64 & 0.11 & 0.89 \\
8 & 256 & 0.37 & 1.02 \\
10 & 1024 & 1.83 & 1.14 \\
12 & 4096 & 8.64 & 1.24 \\
\end{longtable}

\subsection{Comparison with existing tools}\label{sec:comparison}

\textbf{searchnet} occupies a unique niche at the intersection of
fitness landscape simulation and social network modeling.
Table\textasciitilde{}\ref{tab:feature-comparison} summarizes the
feature sets of four relevant tools.

\begin{table}[t]
\centering
\caption{Feature comparison of software tools for fitness landscape and
network-based search simulation.  ``$\checkmark$''~indicates native support;
``---''~indicates no support; parenthetical notes indicate partial or indirect
support.}
\label{tab:feature-comparison}
\begin{tabular}{lcccc}
\toprule
Feature & \textbf{searchnet} & \textbf{nkmodel} & \textbf{RSiena} & NetLogo \\
\midrule
NK landscape generation           & $\checkmark$              & $\checkmark$  & ---           & (ext.)    \\
Multi-actor search                & $\checkmark$ (endogenous) & ---           & $\checkmark$ (networks) & $\checkmark$ (ABM) \\
Bipartite networks                & $\checkmark$              & ---           & $\checkmark$  & (ext.)    \\
Influence matrix ($\mathbf{B}'\mathbf{W}\mathbf{B}$) & $\checkmark$ & (implicit via $K$) & --- & (manual) \\
Statistical estimation            & $\checkmark$ (via RSiena) & ---           & $\checkmark$  & ---       \\
Formal proofs (NK $\subsetneq$ SAOM-NK) & $\checkmark$              & ---           & ---           & ---       \\
Causal inference pipeline         & $\checkmark$ (DID/Synth)  & ---           & ---           & ---       \\
$\{K\}$ framework                    & $\checkmark$              & ---           & ---           & ---       \\
\bottomrule
\end{tabular}
\end{table}

The three most relevant points of comparison are as follows.

\textbf{\textbf{nkmodel}.} The \textbf{nkmodel} package
\citep{Baumann2019nk} provides a faithful implementation of the
classical NK model: it generates random fitness landscapes, computes
fitness values for binary configuration vectors, and supports
hill-climbing search. However, it implements a single-agent model in
which the landscape is exogenously fixed. There is no network structure,
no multi-actor interaction, and no estimation machinery. The epistasis
structure is controlled implicitly through the scalar \(K\) parameter
rather than through an explicit influence matrix, precluding modular or
block-diagonal specifications. \textbf{searchnet} subsumes the
\textbf{nkmodel} use case as a special case: setting \(M = 1\) and
disabling all endogenous network effects recovers the classical NK
model, while the explicit influence matrix \(\mathbf{W}\)
(Section\textasciitilde{}\ref{sec:bipartite}) provides finer-grained
control over interdependence structure.

\textbf{\textbf{RSiena}.} The \textbf{RSiena} package \citep{Ripley2022}
is the reference implementation for stochastic actor-oriented models. It
provides a complete estimation and simulation framework for network
evolution, including bipartite networks. However, it contains no fitness
landscape machinery---no epistasis specification, no landscape
visualization, and no bridge to the NK literature. Using \textbf{RSiena}
alone, a researcher must manually construct the mapping between network
statistics and landscape properties, a task that \textbf{searchnet}
automates through the \(\{K\}\) framework. \textbf{searchnet} uses
\textbf{RSiena} as its simulation engine, wrapping it with a
domain-specific API that translates strategy-theory concepts (epistasis,
fitness, ruggedness) into the corresponding SAOM specification.

\textbf{NetLogo.} NetLogo \citep{Wilensky1999} is widely used for
agent-based simulation in management research. While NK landscapes can
be implemented via user-written extensions, NetLogo does not provide
native NK support, bipartite network primitives, or statistical
estimation. Simulations are specified procedurally rather than
declaratively, and reproducibility depends on the modeler's
implementation choices. \textbf{searchnet}'s declarative specification
(\texttt{saomnk\_model()}) and validation against the classical NK model
provide stronger guardrails for computational experiments.

In summary, \textbf{nkmodel} provides NK landscapes without networks,
\textbf{RSiena} provides networks without landscapes, and NetLogo
provides neither natively. \textbf{searchnet} bridges both traditions,
adding the \(\{K\}\) interpretive framework and a causal inference
pipeline that none of the alternatives offer.

\subsection{Applications}\label{sec:applications}

The \(\{K\}\) formulation is currently exercised by eight independent
research programs. They differ in how much of \textbf{searchnet} they
use, and the distinction matters for a software paper: a package can be
used as a simulation engine, as a post-estimation layer over
\textbf{RSiena}, or as a reference implementation to check an
independent one against. All three modes appear below, and we label each
rather than presenting a uniform façade. Each vignette states a
problematization, the specification that answers it, and a takeaway, so
that a reader can find the closest case to their own problem quickly. We
separate two groups: applications that run \textbf{searchnet} today, and
\emph{refactor candidates} --- studies built on the same \(\{K\}\)
formulation with bespoke code, for which we name the specific capability
that would let the package absorb them. The second group is included
deliberately: the gap between a framework and a package is most visible
where a user would currently have to write their own engine.

{\def\LTcaptype{none} % do not increment counter
\begin{longtable}[]{@{}
  >{\raggedright\arraybackslash}p{(\linewidth - 6\tabcolsep) * \real{0.2500}}
  >{\raggedright\arraybackslash}p{(\linewidth - 6\tabcolsep) * \real{0.2500}}
  >{\raggedright\arraybackslash}p{(\linewidth - 6\tabcolsep) * \real{0.2500}}
  >{\raggedright\arraybackslash}p{(\linewidth - 6\tabcolsep) * \real{0.2500}}@{}}
\toprule\noalign{}
\begin{minipage}[b]{\linewidth}\raggedright
Application
\end{minipage} & \begin{minipage}[b]{\linewidth}\raggedright
Domain
\end{minipage} & \begin{minipage}[b]{\linewidth}\raggedright
Mode of use
\end{minipage} & \begin{minipage}[b]{\linewidth}\raggedright
\(\{K\}\) channels driven
\end{minipage} \\
\midrule\noalign{}
\endhead
\bottomrule\noalign{}
\endlastfoot
Competitive convergence & U.S. airlines, 1993--2025 & Reference
implementation, conformance-checked & \(K_{AC}\), \(K_{CC}\),
\(K_{CA}\), \(K_{AA}\) \\
Architectural shocks & Repositioning under innovation & Simulation
engine & \(K_{CC}\), \(K_{AA}\), rate \\
Workforce formation & Startup hiring & Simulation engine, calibrated &
\(K_{AC}\), \(K_{CA}\), \(K_{CC}\) \\
Policy shocks & Subsidies and risk-taking & Simulation engine, shock
schedule & \(K_{AC}\), \(K_{CC}\), covariate \\
Firm vs.~industry effects & Variance decomposition & Simulation engine &
\(K_{AC}\), \(K_{CA}\), \(K_{CC}\) \\
Market-access institutions & Cross-border expansion & Post-estimation
layer & \(K_{AC}\), \(K_{CA}\), \(K_{CC}\) \\
Evaluator discipline & Venture capital & Bespoke code; time-varying
\(W\) added in 0.7.0 & \(K_{AC}\), \(K_{CA}\), \(K_{CC}\) \\
Basin geometry & Industrial policy & Refactor candidate (tournament
driver) & \(K_{AC}\), \(K_{CC}\), \(K_{AA}\) \\
\end{longtable}
}

Two channels deserve a warning that applies throughout. First,
\(K_{CA}\) is driven by \texttt{inPop} (degree popularity), not by
similarity-weighted imitation. \textbf{RSiena}'s effect set for a
bipartite dependent variable contains no similarity effect:
\texttt{simEgoInDist2} is defined for one-mode networks only.
\textbf{searchnet} supplies the statistic for post-hoc computation but
cannot make it a simulable evaluation term, and no application below
uses it as one. Second, \(K_{AA}\) is a \emph{projection}, so it is
frequently measured as an outcome rather than entered as an effect; the
table records what is driven, not what is reported.

\subsubsection{Competitive convergence in
U.S.~airlines}\label{sec:app-airlines}

\textbf{Research question.} Rival firms often come to resemble one
another, and the resemblance is routinely read as evidence of imitation.
Does it license that inference? Between 1993 and 2025 pairwise route
overlap among U.S.~carriers rose 4.8-fold (mean pairwise Jaccard
\(0.0076 \rightarrow 0.0364\)) while concentration on contested routes
\emph{declined} --- a pattern imitation accounts do not predict.

\textbf{Construct and channels.} The claim under test is
\emph{architectural channeling}: that shared activity-system structure
alone funnels independent, boundedly rational search into overlapping
configurations, so convergence arises with no firm copying another. This
is a statement about \(K_{CC}\) (the interaction matrix \(W\)) producing
\(K_{AA}\) (pairwise overlap) without \(K_{CA}\) (imitation) doing any
work, and it is testable only counterfactually --- by deleting channels
one at a time and re-running the same seeds.

\textbf{Design.} \(M = 30\) firms, \(N = 256\) activities, 100 rounds;
\(W\) parameterised by pattern (random, modular, hierarchical) crossed
with epistatic density; the main sweep crosses 2 densities \(\times\) 8
rivalry \(\times\) 8 imitation values, 10 replications each, for
1\{,\}280 runs.

\textbf{Results.} Deleting channels from the calibrated model separates
them cleanly. Terminal \(\bar K_{AA}\) at round 100, ten replications
per scenario:

{\def\LTcaptype{none} % do not increment counter
\begin{longtable}[]{@{}llll@{}}
\toprule\noalign{}
Scenario & \(\bar K_{AA}\) & 95\% interval & Fold \\
\midrule\noalign{}
\endhead
\bottomrule\noalign{}
\endlastfoot
Baseline (calibrated) & 0.034 & 0.020--0.052 & 1.00 \\
Imitation removed (\(K_{CA} = 0\)) & 0.108 & 0.086--0.128 &
\textbf{3.2} \\
Epistasis removed (\(K_{CC} = 0\)) & 0.010 & 0.007--0.017 & 0.29 \\
All competition removed & 0.399 & 0.322--0.543 & 11.7 \\
\end{longtable}
}

Removing imitation \emph{raises} convergence 3.2-fold, and removing
epistasis collapses it by 71 percent. Under similarity-weighted
imitation firms chase whichever peers resemble them, which disperses
them across the space; shared structure is what concentrates them. A
54-cell robustness sweep finds the ordering holds in 27 of 27
hierarchical cells and 17 of 27 modular cells, which scopes the claim to
architectural regime rather than asserting it universally.

\begin{figure}
\centering
\includegraphics[width=1\linewidth,height=\textheight,keepaspectratio,alt={Industry trajectory. Mean pairwise route overlap among U.S. carriers 1993--2025 with era shading and merger markers, the empirical series the simulation is calibrated against.}]{searchnet-jss-NEW2_files/figures/cd2026_kaa_trajectory.png}
\caption{Industry trajectory. Mean pairwise route overlap among U.S.
carriers 1993--2025 with era shading and merger markers, the empirical
series the simulation is calibrated against.\label{fig:kaatraj}}
\end{figure}

\textbf{Mode of use, stated plainly.} This application predates the
package. Its simulation results come from an independent procedural
implementation, and \textbf{searchnet} enters as a second implementation
checked against it. That arrangement is worth describing because it is
how the package is validated in practice. A conformance harness runs
nine parameter cells \(\times\) three seeds through both implementations
and compares them on three layers: structural agreement, trajectory
hashes, and \emph{invariants} --- the NK reduction of Theorem 1
(Reduction) and monotonicity of crowding in the rivalry parameter. Its
design note states the discipline that makes such a comparison
meaningful:

\begin{quote}
\textbf{searchnet} is R6 over \textbf{RSiena} and the reference engine
is procedural, so bit-identical trajectories are \emph{not} expected:
the random number stream is consumed in a different order. Reporting a
cross-paradigm trajectory mismatch as a defect would be a category
error.
\end{quote}

Structural and invariant layers must agree; trajectory hashes need not.
We recommend this three-layer split to anyone validating a simulation
package against an existing implementation.

A second discipline this application enforces is worth transferring. The
potential-game property of
Section\textasciitilde{}\ref{sec:formal-properties} was tested \emph{on
the simulation code} rather than on a reduced model: add an activity,
then drop it, and net utility must return to zero. It does, to machine
precision, for the core specification. The same probe \emph{locates} the
boundary --- with similarity-weighted imitation active the path integral
deviates by 2.68 at calibrated coefficients, because that term depends
on rivals' performances, which move when the focal firm moves. The
formal guarantees are therefore claimed for the core specification only,
and the probe is what draws the line.

\textbf{Takeaway.} Convergence is produced by shared structure, not by
copying: the channel whose removal \emph{raises} overlap cannot be the
channel generating it. For package users, the transferable practice is
the three-layer conformance split and the path-integral probe that marks
where a formal guarantee stops.

\subsubsection{Architectural shocks and the rate--evaluation
split}\label{sec:app-shocks}

\textbf{Research question.} When an innovation reshapes which activities
complement which, firms differ in how fast they can respond. Can the
\emph{cost of effecting} a repositioning be separated from the
\emph{value of the destination} using observational data? Archival
panels record the option taken and never the options declined, so ``did
not reposition'' is equally consistent with \emph{could not afford to}
and \emph{did not want to}.

\textbf{Construct and channels.} This is the application that exploits
the SAOM's distinctive factorisation. The rate function \(\lambda_i\)
governs \emph{how often} an actor gets to move; the evaluation function
\(f_i\) governs \emph{where} it moves. Adjustment cost is a rate
quantity and destination value is an evaluation quantity, so a model
that separates them can identify what a single-equation model cannot.
The design manipulates \texttt{RateX} alone and holds the evaluation
function \emph{identical} across arms --- the covariate effect is
declared with \texttt{parameter = 0} precisely so that it is inert.

\textbf{Design.} \(M = 12\) actors, \(N = 12\) components, 30 iterations
per actor (360 ministeps), 20 seed pairs. \(K_{CC}\) shifts mid-run: two
\texttt{XWX} dyadic covariates carry \(W_{\text{pre}}\) and
\(W_{\text{post}}\), and \texttt{theta\_shocks} flips their weights from
\((1,0)\) to \((0,1)\) at the midpoint, so the firm's position carries
across the switch and is re-scored under the new structure.

\begin{figure}
\centering
\includegraphics[width=1\linewidth,height=\textheight,keepaspectratio,alt={The K\_\{CC\} switch. Left, the pre-shock influence matrix W\_\{\textbackslash text\{pre\}\}; centre, the post-shock W\_\{\textbackslash text\{post\}\}; right, their difference. The shock schedule flips the weights of two  covariates at the midpoint, so actors carry their positions across a change in which components complement which.}]{searchnet-jss-NEW2_files/figures/ishock_w_switch.png}
\caption{The \(K_{CC}\) switch. Left, the pre-shock influence matrix
\(W_{\text{pre}}\); centre, the post-shock \(W_{\text{post}}\); right,
their difference. The shock schedule flips the weights of two
\texttt{XWX} covariates at the midpoint, so actors carry their positions
across a change in which components complement
which.\label{fig:wswitch}}
\end{figure}

\begin{Shaded}
\begin{Highlighting}[]
\NormalTok{Wpre  }\OtherTok{\textless{}{-}} \FunctionTok{create\_block\_diag}\NormalTok{(N, }\DecValTok{2}\NormalTok{)}
\NormalTok{Wpost }\OtherTok{\textless{}{-}} \FunctionTok{create\_block\_diag}\NormalTok{(N, }\DecValTok{3}\NormalTok{)}

\NormalTok{mkm }\OtherTok{\textless{}{-}} \ControlFlowTok{function}\NormalTok{(rateX) \{}
\NormalTok{  rts }\OtherTok{\textless{}{-}} \ControlFlowTok{if}\NormalTok{ (rateX }\SpecialCharTok{==} \DecValTok{0}\NormalTok{) }\FunctionTok{list}\NormalTok{() }\ControlFlowTok{else} \FunctionTok{list}\NormalTok{(}
    \FunctionTok{list}\NormalTok{(}\AttributeTok{effect =} \StringTok{"RateX"}\NormalTok{, }\AttributeTok{parameter =}\NormalTok{ rateX, }\AttributeTok{dv\_name =}\NormalTok{ DV, }\AttributeTok{fix =} \ConstantTok{TRUE}\NormalTok{,}
         \AttributeTok{interaction1 =} \StringTok{"self$strat\_1\_coCovar"}\NormalTok{))}
  \FunctionTok{list}\NormalTok{(}\AttributeTok{dv\_bipartite =} \FunctionTok{list}\NormalTok{(}\AttributeTok{name =}\NormalTok{ DV, }\AttributeTok{rates =}\NormalTok{ rts,}
    \AttributeTok{effects =} \FunctionTok{list}\NormalTok{(}
      \FunctionTok{list}\NormalTok{(}\AttributeTok{effect =} \StringTok{"density"}\NormalTok{, }\AttributeTok{parameter =} \SpecialCharTok{{-}}\FloatTok{0.6}\NormalTok{, }\AttributeTok{dv\_name =}\NormalTok{ DV, }\AttributeTok{fix =} \ConstantTok{TRUE}\NormalTok{),}
      \FunctionTok{list}\NormalTok{(}\AttributeTok{effect =} \StringTok{"inPop"}\NormalTok{,   }\AttributeTok{parameter =} \SpecialCharTok{{-}}\FloatTok{0.4}\NormalTok{, }\AttributeTok{dv\_name =}\NormalTok{ DV, }\AttributeTok{fix =} \ConstantTok{TRUE}\NormalTok{)),}
    \AttributeTok{coCovars =} \FunctionTok{list}\NormalTok{(}
      \FunctionTok{list}\NormalTok{(}\AttributeTok{effect =} \StringTok{"egoX"}\NormalTok{, }\AttributeTok{parameter =} \DecValTok{0}\NormalTok{, }\AttributeTok{dv\_name =}\NormalTok{ DV, }\AttributeTok{fix =} \ConstantTok{TRUE}\NormalTok{,}
           \AttributeTok{interaction1 =} \StringTok{"self$strat\_1\_coCovar"}\NormalTok{, }\AttributeTok{x =}\NormalTok{ grp)),}
    \AttributeTok{coDyadCovars =} \FunctionTok{list}\NormalTok{(}
      \FunctionTok{list}\NormalTok{(}\AttributeTok{effect =} \StringTok{"XWX"}\NormalTok{, }\AttributeTok{parameter =} \DecValTok{1}\NormalTok{, }\AttributeTok{dv\_name =}\NormalTok{ DV, }\AttributeTok{fix =} \ConstantTok{TRUE}\NormalTok{,}
           \AttributeTok{nodeSet =} \FunctionTok{c}\NormalTok{(}\StringTok{"COMPONENTS"}\NormalTok{, }\StringTok{"COMPONENTS"}\NormalTok{),}
           \AttributeTok{interaction1 =} \StringTok{"self$component\_1\_coDyadCovar"}\NormalTok{, }\AttributeTok{x =}\NormalTok{ Wpre),}
      \FunctionTok{list}\NormalTok{(}\AttributeTok{effect =} \StringTok{"XWX"}\NormalTok{, }\AttributeTok{parameter =} \DecValTok{0}\NormalTok{, }\AttributeTok{dv\_name =}\NormalTok{ DV, }\AttributeTok{fix =} \ConstantTok{TRUE}\NormalTok{,}
           \AttributeTok{nodeSet =} \FunctionTok{c}\NormalTok{(}\StringTok{"COMPONENTS"}\NormalTok{, }\StringTok{"COMPONENTS"}\NormalTok{),}
           \AttributeTok{interaction1 =} \StringTok{"self$component\_2\_coDyadCovar"}\NormalTok{, }\AttributeTok{x =}\NormalTok{ Wpost))))}
\NormalTok{\}}

\NormalTok{shocks }\OtherTok{\textless{}{-}} \FunctionTok{list}\NormalTok{(}
  \FunctionTok{list}\NormalTok{(}\AttributeTok{effect\_level =} \FunctionTok{c}\NormalTok{(}\StringTok{"XWX\_1"}\NormalTok{, }\StringTok{"XWX\_2"}\NormalTok{), }\AttributeTok{parameter =} \FunctionTok{c}\NormalTok{(}\DecValTok{1}\NormalTok{, }\DecValTok{0}\NormalTok{), }\AttributeTok{portion =} \DecValTok{1}\NormalTok{, }\AttributeTok{shock\_on =} \DecValTok{0}\NormalTok{),}
  \FunctionTok{list}\NormalTok{(}\AttributeTok{effect\_level =} \FunctionTok{c}\NormalTok{(}\StringTok{"XWX\_1"}\NormalTok{, }\StringTok{"XWX\_2"}\NormalTok{), }\AttributeTok{parameter =} \FunctionTok{c}\NormalTok{(}\DecValTok{0}\NormalTok{, }\DecValTok{1}\NormalTok{), }\AttributeTok{portion =} \DecValTok{1}\NormalTok{, }\AttributeTok{shock\_on =} \DecValTok{1}\NormalTok{))}

\NormalTok{env }\OtherTok{\textless{}{-}}\NormalTok{ SaomNkRSienaBiEnv}\SpecialCharTok{$}\FunctionTok{new}\NormalTok{(}\FunctionTok{list}\NormalTok{(}\AttributeTok{M =}\NormalTok{ M, }\AttributeTok{N =}\NormalTok{ N, }\AttributeTok{BI\_PROB =} \FloatTok{0.3}\NormalTok{, }\AttributeTok{rand\_seed =} \DecValTok{7000} \SpecialCharTok{+}\NormalTok{ s))}
\NormalTok{env}\SpecialCharTok{$}\FunctionTok{search\_rsiena}\NormalTok{(}\AttributeTok{structure\_model =} \FunctionTok{mkm}\NormalTok{(rateX), }\AttributeTok{iterations\_per\_actor =} \DecValTok{30}\NormalTok{,}
                  \AttributeTok{run\_seed =} \DecValTok{7000} \SpecialCharTok{+}\NormalTok{ s, }\AttributeTok{theta\_shocks =}\NormalTok{ shocks)}
\end{Highlighting}
\end{Shaded}

\textbf{Results.} The rate manipulation separates cleanly: with
\texttt{RateX = 2}, quick-to-act actors make 51.0 tie changes against
6.7 for slow-to-act, a ratio of 7.6, while the control arm sits at 1.05.
Every one of 40 runs shows a post-shock fitness valley. Two results are
reported as null and should be: neither valley depth (\(p = 0.18\)) nor
time to trough (\(p = 0.50\)) differs between arms.

The substantive finding is a selection wedge. Comparing movers to
non-movers in the same treated world gives \(-0.054\); comparing each
mover to its own matched counterfactual gives \(-0.032\); the
difference, \(-0.022\) (\(t = -2.46\), \(p = 0.015\)), is what the
observational comparison attributes to repositioning that is in fact
selection. Notably, applying the \emph{airline} project's decomposition
to the same runs gives a mean wedge of \(-0.002\) --- degenerate,
exactly as that decomposition's algebraic identity requires. Two
conventions, one dataset, opposite informativeness: a reminder that a
decomposition's benchmark determines what it can detect.

\begin{figure}
\centering
\includegraphics[width=1\linewidth,height=\textheight,keepaspectratio,alt={Rate separation. With  active, quick-to-act actors make 51.0 tie changes against 6.7 for slow-to-act (ratio 7.6) while the control arm sits at 1.05, with evaluation functions held identical across arms.}]{searchnet-jss-NEW2_files/figures/ishock_rate_separation.png}
\caption{Rate separation. With \texttt{RateX} active, quick-to-act
actors make 51.0 tie changes against 6.7 for slow-to-act (ratio 7.6)
while the control arm sits at 1.05, with evaluation functions held
identical across arms.\label{fig:rateseparation}}
\end{figure}

\textbf{Takeaway.} Separating the rate function from the evaluation
function turns an unidentified question --- was the firm unwilling or
unable? --- into two estimable quantities. A decomposition's benchmark
determines what it can detect: the same runs yield a significant wedge
under one convention and a degenerate zero under another.

\subsubsection{Startup workforce formation}\label{sec:app-workforce}

\textbf{Research question.} Startups hire to build capability, but
breadth carries a coordination cost. When does successful
co-construction foreclose its own continuation? The Penrose Effect
predicts an inverted-U between workforce breadth and performance.

\textbf{Construct and channels.} Actors are startups, components are
skills. Workforce breadth is \(K_{AC}\) (\texttt{outAct}); competition
for the same scarce skills is \(K_{CA}\) (\texttt{inPop}); skill
complementarity within a team is \(K_{CC}\), entered as an empirically
estimated \(W\). Labour-market overlap between startups is \(K_{AA}\)
and is \emph{measured but not driven} --- a good illustration of the
projection channel used as an outcome.

\textbf{Design and calibration.} \(M = 50\) startups, \(N = 12\) skills,
20 steps per actor. The parameters are not chosen; they are
\textbf{RSiena} point estimates from a fitted bipartite SAOM on
Crunchbase founding cohorts (\texttt{density} \(= -1.756\),
\texttt{inPop} \(= 0.022\), \texttt{outAct} \(= 0.089\)), carried into
simulation through the functional API:

\begin{Shaded}
\begin{Highlighting}[]
\NormalTok{env }\OtherTok{\textless{}{-}} \FunctionTok{saomnk\_env}\NormalTok{(}\AttributeTok{M =} \DecValTok{50}\NormalTok{, }\AttributeTok{N =} \DecValTok{12}\NormalTok{, }\AttributeTok{density =} \FloatTok{0.3}\NormalTok{, }\AttributeTok{seed =}\NormalTok{ seed)}

\NormalTok{mod }\OtherTok{\textless{}{-}} \FunctionTok{saomnk\_model}\NormalTok{(}
  \AttributeTok{density          =} \SpecialCharTok{{-}}\FloatTok{1.756}\NormalTok{,   }\CommentTok{\# fitted}
  \AttributeTok{popularity       =}  \FloatTok{0.022}\NormalTok{,   }\CommentTok{\# fitted (K\_CA)}
  \AttributeTok{scope            =}  \FloatTok{0.089}\NormalTok{,   }\CommentTok{\# fitted (K\_AC)}
  \AttributeTok{influence\_matrix =}\NormalTok{ W\_emp,    }\CommentTok{\# empirical influence matrix (drives K\_CC)}
  \AttributeTok{influence\_weight =} \SpecialCharTok{{-}}\FloatTok{0.062}\NormalTok{)}

\FunctionTok{saomnk\_run}\NormalTok{(env, mod, }\AttributeTok{steps\_per\_actor =} \DecValTok{20}\DataTypeTok{L}\NormalTok{, }\AttributeTok{seed =}\NormalTok{ seed }\SpecialCharTok{+} \DecValTok{7}\DataTypeTok{L}\NormalTok{)}
\NormalTok{degs }\OtherTok{\textless{}{-}} \FunctionTok{saomnk\_get\_degrees}\NormalTok{(env)   }\CommentTok{\# K\_AC, K\_CA, K\_AA, K\_CC read out together}
\end{Highlighting}
\end{Shaded}

\begin{figure}
\centering
\includegraphics[width=0.85\linewidth,height=\textheight,keepaspectratio,alt={Utility decomposition for the workforce application, separating each actor's objective into its density, scope and epistasis contributions. The scope term is positive at the fitted parameters, which is why no inverted U appears.}]{searchnet-jss-NEW2_files/figures/bort_utility_waterfall.png}
\caption{Utility decomposition for the workforce application, separating
each actor's objective into its density, scope and epistasis
contributions. The scope term is positive at the fitted parameters,
which is why no inverted U appears.\label{fig:bortwaterfall}}
\end{figure}

\textbf{Results, including a negative one.} Under the calibrated
specification the inverted-U \emph{does not appear}: the quadratic
coefficient is positive in all seven conditions. The diagnosis is
transparent rather than mysterious --- the fitted \texttt{outAct} is
\(+0.089\), so the estimated model contains no scope-cost term, and a
curve with no cost cannot bend down. Re-specifying with \texttt{scope}
\(< 0\) recovers the inverted-U in 8 of 27 cells, and every one of those
8 sits at the empirical popularity value. A further design lets the
scope cost regenerate from the realised landscape each wave; the
inverted-U appears, but the emergence ratio is \(\approx 0.99\), meaning
the curvature passes through from the imposed cost rather than being
generated by feedback. We report this because it is the kind of result a
simulation package should make easy to obtain and hard to overstate: the
mechanism was not found, and the reason it was not found is legible in
the fitted parameters.

\includegraphics[width=0.49\linewidth,height=\textheight,keepaspectratio,alt={Workforce application. Left, the \textbackslash\{K\textbackslash\}-4 panel for 50 startups over 250 steps, showing all four channels read out together. Right, breadth (K\_\{AC\}) against utility; the relationship is monotone and concave, that is diminishing returns, not the inverted U the Penrose Effect predicts.}]{searchnet-jss-NEW2_files/figures/bort_k4_panel.png}
\includegraphics[width=0.49\linewidth,height=\textheight,keepaspectratio]{searchnet-jss-NEW2_files/figures/bort_phase_space.png}

\textbf{Takeaway.} Calibrating to fitted coefficients rather than to
chosen ones makes a simulation falsifiable, and here it falsified the
expected result: no inverted U appears because the fitted scope term is
positive, so the model contains no coordination cost to bend the curve.

\subsubsection{Policy shocks and risk-taking}\label{sec:app-policy}

\textbf{Research question.} When a government subsidy arrives and later
withdraws, do treated firms shift toward exploratory activity, and does
the shift persist once the support is removed?

\textbf{Construct and channels.} A subsidy is modelled as a timed change
in an actor-covariate effect rather than as a change in the landscape:
eligible firms receive a positive \texttt{egoX} weight, together with
interactions on scope (\texttt{outAct|egoX}) that make the subsidy's
effect depend on how broad the firm already is. \(K_{CC}\) carries the
exploration structure through \texttt{XWX} with a component risk matrix.
The design exercises \textbf{searchnet}'s shock schedule in its
three-regime form --- pre-policy, policy, post-policy --- which is the
natural way to represent a temporary intervention:

\begin{Shaded}
\begin{Highlighting}[]
\NormalTok{theta\_shocks }\OtherTok{\textless{}{-}} \FunctionTok{list}\NormalTok{(}
  \FunctionTok{list}\NormalTok{(}\AttributeTok{effect\_level =} \FunctionTok{c}\NormalTok{(}\StringTok{"egoX"}\NormalTok{, }\StringTok{"outAct|egoX\_1"}\NormalTok{, }\StringTok{"outActSqrt|egoX\_2"}\NormalTok{),}
       \AttributeTok{parameter =} \FunctionTok{c}\NormalTok{(}\DecValTok{0}\NormalTok{, }\DecValTok{0}\NormalTok{, }\DecValTok{0}\NormalTok{),             }\AttributeTok{portion =}\NormalTok{ pre\_portion,    }\AttributeTok{shock\_on =} \DecValTok{0}\NormalTok{),}
  \FunctionTok{list}\NormalTok{(}\AttributeTok{effect\_level =} \FunctionTok{c}\NormalTok{(}\StringTok{"egoX"}\NormalTok{, }\StringTok{"outAct|egoX\_1"}\NormalTok{, }\StringTok{"outActSqrt|egoX\_2"}\NormalTok{),}
       \AttributeTok{parameter =} \FunctionTok{c}\NormalTok{(}\DecValTok{1}\NormalTok{, h1\_val, h2\_val),   }\AttributeTok{portion =}\NormalTok{ policy\_portion, }\AttributeTok{shock\_on =} \DecValTok{1}\NormalTok{),}
  \FunctionTok{list}\NormalTok{(}\AttributeTok{effect\_level =} \FunctionTok{c}\NormalTok{(}\StringTok{"egoX"}\NormalTok{, }\StringTok{"outAct|egoX\_1"}\NormalTok{, }\StringTok{"outActSqrt|egoX\_2"}\NormalTok{),}
       \AttributeTok{parameter =} \FunctionTok{c}\NormalTok{(}\DecValTok{0}\NormalTok{, }\DecValTok{0}\NormalTok{, }\DecValTok{0}\NormalTok{),             }\AttributeTok{portion =}\NormalTok{ post\_portion,   }\AttributeTok{shock\_on =} \DecValTok{0}\NormalTok{))}

\NormalTok{env }\OtherTok{\textless{}{-}}\NormalTok{ SaoMNK}\SpecialCharTok{$}\FunctionTok{new}\NormalTok{(}\FunctionTok{list}\NormalTok{(}\AttributeTok{M =}\NormalTok{ M, }\AttributeTok{N =}\NormalTok{ N, }\AttributeTok{init\_matrix =}\NormalTok{ init\_matrix))}
\NormalTok{env}\SpecialCharTok{$}\FunctionTok{search\_rsiena}\NormalTok{(}\AttributeTok{structure\_model =}\NormalTok{ structure\_model,}
                  \AttributeTok{theta\_shocks    =}\NormalTok{ theta\_shocks,}
                  \AttributeTok{iterations\_per\_actor =}\NormalTok{ n\_iterations\_per\_actor }\SpecialCharTok{*}\NormalTok{ env}\SpecialCharTok{$}\NormalTok{N,}
                  \AttributeTok{run\_seed =}\NormalTok{ run\_seed)}
\end{Highlighting}
\end{Shaded}

\textbf{Design.} \(M = 18\) firms, \(N = 16\) components, 64 iterations
per actor; a \(6 \times 5 \times 4 \times 3\) parameter grid of 360
cells at five seeds each, for 1\{,\}800 simulations.

\textbf{What the package supplies.} The risk-taking trajectory and
exploration--exploitation plots are produced by \textbf{searchnet}'s own
plotting layer from the simulated chain, not by bespoke code, which is
the intended division of labour: the package computes the actor-level
series and the analyst supplies the contrast. We show the trajectory
panel rather than the difference-in-differences panel, because the
latter's period axis is ordered lexically rather than temporally --- a
reminder that a plot returned by a package is still the analyst's
responsibility to read.

\begin{figure}
\centering
\includegraphics[width=1\linewidth,height=\textheight,keepaspectratio,alt={Policy shock. Actor-level risk-taking trajectories for subsidised and control firms across the three-regime shock schedule, plotted by  from the simulated chain; the vertical line marks the policy onset.}]{searchnet-jss-NEW2_files/figures/amr_risk_trajectories.png}
\caption{Policy shock. Actor-level risk-taking trajectories for
subsidised and control firms across the three-regime shock schedule,
plotted by \textbf{searchnet} from the simulated chain; the vertical
line marks the policy onset.\label{fig:policyshock}}
\end{figure}

\textbf{Takeaway.} A temporary intervention is naturally expressed as a
three-regime shock schedule, and the package's plotting layer returns
actor-level series ready for the contrast the analyst chooses.

\subsubsection{Firm and industry effects}\label{sec:app-variance}

The variance-decomposition literature asks whether firm or industry
effects explain more of the variance in profitability. If firms are
competitively coupled, the two are not separable in principle, and a
simulation in which the data-generating process is known can show what a
decomposition recovers.

This application is the cleanest illustration of the minimal functional
pipeline, and we present it as such:

\begin{Shaded}
\begin{Highlighting}[]
\NormalTok{env }\OtherTok{\textless{}{-}} \FunctionTok{saomnk\_env}\NormalTok{(}\AttributeTok{M =}\NormalTok{ M, }\AttributeTok{N =}\NormalTok{ N, }\AttributeTok{density =} \FloatTok{0.2}\NormalTok{, }\AttributeTok{seed =}\NormalTok{ seed)}
\NormalTok{K   }\OtherTok{\textless{}{-}} \FunctionTok{make\_influence\_matrix}\NormalTok{(N, arch\_type)      }\CommentTok{\# modular | hierarchical | random}

\NormalTok{mod }\OtherTok{\textless{}{-}} \FunctionTok{saomnk\_model}\NormalTok{(}\AttributeTok{density =} \SpecialCharTok{{-}}\FloatTok{0.5}\NormalTok{, }\AttributeTok{popularity =}\NormalTok{ coupling,}
                    \AttributeTok{scope =} \FloatTok{0.05}\NormalTok{, }\AttributeTok{influence\_matrix =}\NormalTok{ K, }\AttributeTok{influence\_weight =} \FloatTok{0.1}\NormalTok{)}

\FunctionTok{saomnk\_run}\NormalTok{(env, mod, }\AttributeTok{steps\_per\_actor =} \DecValTok{15}\DataTypeTok{L}\NormalTok{, }\AttributeTok{seed =}\NormalTok{ seed }\SpecialCharTok{+} \DecValTok{1}\DataTypeTok{L}\NormalTok{)}
\NormalTok{bi\_mat }\OtherTok{\textless{}{-}} \FunctionTok{saomnk\_get\_bipartite}\NormalTok{(env)}
\end{Highlighting}
\end{Shaded}

Coupling strength enters through \texttt{popularity}, sweeping
\(\{0, 0.1, 0.3,
0.5, 0.8\}\) across three architectures. We report the pipeline, not a
bias estimate: in the present implementation the ``true'' effects are
re-derived from the same realised performance the estimator sees, so the
recovered bias is zero by construction. Establishing a bias requires
taking the truth from the generating parameters instead, and that work
is ongoing. We state this because a software paper is the wrong place to
launder an unfinished result, and because the pipeline above is
genuinely the shortest path from a specification to a simulated
bipartite panel.

\textbf{Takeaway.} The shortest path from specification to simulated
panel is four calls. The vignette also shows what a simulation cannot do
for you: if the ``true'' effects are re-derived from the same realised
outcome the estimator sees, recovered bias is zero by construction,
whatever the coupling.

\subsubsection{Market-access institutions}\label{sec:app-international}

\textbf{Research question.} Do institutional distances between countries
act as tolls on entry, or as topology that reshapes which markets are
complementary? A toll shifts the cost of a market; topology changes what
a market is worth \emph{given what else the firm holds} --- and the two
are separable, because the first is a main effect on the component and
the second is \(K_{CC}\).

\textbf{Design.} A bipartite firm-by-market panel: 534 firms \(\times\)
63 markets \(\times\) 4 waves, 8\{,\}023 ties, with nine coupling
matrices built from CEPII gravity data and United Nations
voting-similarity scores. Each matrix enters \emph{twice} --- once as a
component main effect and once as \texttt{XWX} --- so the two mechanisms
compete for the same variance in one model.

\includegraphics[width=0.49\linewidth,height=\textheight,keepaspectratio,alt={Post-estimation use. Left, a geographic network view produced by . Right, a specification curve built with 's estimate-extraction and scaled-autocorrelation helpers over  fits; these functions existed but were unreachable without  until the  was corrected in 0.6.0.}]{searchnet-jss-NEW2_files/figures/intl_geo_network.png}
\includegraphics[width=0.49\linewidth,height=\textheight,keepaspectratio]{searchnet-jss-NEW2_files/figures/intl_spec_curve.png}

\textbf{Mode of use and result.} Estimation runs through \textbf{RSiena}
directly; \textbf{searchnet} is used for the specification curve and
network visualisation. Both channels operate: a block Wald test gives
\(\chi^2 = 169.4\) for coupling against \(143.9\) for the main effects,
both on 10 degrees of freedom. An earlier and stronger claim --- that
institutions are topology \emph{rather than} toll --- was withdrawn
after review found the cost-shifting term had been operationalised as a
country-constant row mean while coupling was firm-varying, which is not
a fair contest. The surviving claim is the weaker and better supported
one, that the distance literature is incomplete rather than
misspecified.

\textbf{Takeaway.} Entering each matrix twice, as a component main
effect and as a coupling term, lets two rival mechanisms compete for the
same variance in one model instead of being compared across models.

\subsubsection{Refactor candidates}\label{sec:app-refactor}

Two further studies are built on the \(\{K\}\) formulation but implement
it in bespoke code. We include them because the capability each would
need is specific and short, and naming it is more useful to a
prospective user than omitting the case.

\textbf{Evaluator discipline under capital abundance.}
\emph{Problematization:} when capital becomes abundant, which margins of
organizational evaluation adjust, and what insulates the rest? Investors
are actors and technology categories are components, so an investor's
portfolio is a bipartite row and complementarity between categories is
exactly \(K_{CC}\). \emph{Methods:} a bipartite SAOM on 368 investors
\(\times\) 45 categories over six waves, with \texttt{outActSqrt},
\texttt{inPopSqrt} and \texttt{cycle4}, plus complementarity entering as
a \emph{time-varying actor-by-component} covariate interacted with a
policy-regime indicator. The regime interaction is strongly negative,
indicating that complementarity weighs less in commitment decisions when
capital is cheap. \emph{Takeaway:} the substantive finding is a
\(K_{CC}\) result, but it is not currently expressed as one.
\emph{Capability required:} the study needs a \textbf{time-varying
\(W\)} --- \texttt{varDyadCovar} rather than \texttt{coDyadCovar} --- so
that epistasis can shift across waves. Version 0.7.0 supplies this as
the \texttt{influence\_arrays} argument to \texttt{saomnk\_model},
covering specification and simulation; estimation and visualisation of a
time-varying \(W\) remain open. The analysis reported here predates that
release and uses hand-built covariate arrays.

\textbf{Basin geometry under policy regimes.} \emph{Problematization:}
policy archetypes that produce indistinguishable average fitness may
still leave firms in very differently shaped basins, and it is the
shape, not the height, that determines what a shock does.
\emph{Methods:} 20 firms and 16 activities across 27 conditions crossing
epistatic density, interaction pattern and erosion, with a ten-term
utility whose scope, synergy, herding and congestion terms are \(\{K\}\)
channels written out longhand, and \textbf{RSiena} estimates supplying
the coefficients. \emph{Takeaway:} archetypes converge on fitness --- a
spread of 0.18 percent across four archetypes --- while diverging by
roughly an order of magnitude on hysteresis area. Fitness parity
conceals geometric divergence, which is an argument for reporting basin
shape alongside performance. \emph{Capability required:}
\textbf{searchnet} already computes fitness landscapes and basin
statistics (Section\textasciitilde{}\ref{sec:basins}); what is missing
is a \textbf{tournament driver} that runs heterogeneous strategy
populations against one another and returns ranked outcomes. With that,
the ten-term utility collapses into a \texttt{saomnk\_model} call plus a
scoring function.

We flag one caution these two cases share. Both were written before the
package matured, and both re-implement channels the package now
provides. A reader adopting \textbf{searchnet} should expect the reverse
migration to be the common case: the framework is easier to express in
the package than to re-derive, and the value of doing so is that the
\(\{K\}\) channels then carry their standard meanings rather than local
ones.

\subsubsection{\texorpdfstring{Failure modes inherited from
\textbf{RSiena}}{Failure modes inherited from }}\label{sec:app-failures}

Because \textbf{searchnet} delegates estimation to \texttt{siena07}, it
inherits that estimator's failure modes, and users should know two of
them.

The first is provable. Under method-of-moments estimation, no effect
whose observed statistic lies exactly on the boundary of its attainable
support can converge: the estimate diverges as \(\theta \to \pm\infty\).
In a \emph{balanced} bipartite panel every actor holds at least one
component in every wave, so the isolate count is exactly zero and
\texttt{outTrunc(1)} is exactly \(M\) --- both saturated. A
demonstration on public data changes only the sampling rule: on a
balanced construction \texttt{outIso} estimates \(-18.65\) (SE
\(6.62\)), a runaway; on an unbalanced construction of the same network,
where 1\{,\}547 of 1\{,\}316 actor-waves are isolates, it estimates
\(+3.29\) (SE \(0.42\)). The screen is one line --- compute
\((\text{observed} - \text{lo}) / (\text{hi} - \text{lo})\) and refuse
exactly \(0\) or \(1\) --- and it costs nothing to run before
estimating. We note that a stronger version of this claim, covering
\emph{near}-boundary statistics, was tested and refuted: only exact
saturation disqualifies.

\begin{figure}
\centering
\includegraphics[width=1\linewidth,height=\textheight,keepaspectratio,alt={Why an exactly saturated statistic cannot converge. In a balanced panel the observed isolate count sits on the boundary of its support, so method of moments can match it only as \textbackslash theta \textbackslash to -\textbackslash infty; the same network sampled to permit entry and exit puts the statistic in the interior and the effect estimates cleanly.}]{searchnet-jss-NEW2_files/figures/remedies_boundary.png}
\caption{Why an exactly saturated statistic cannot converge. In a
balanced panel the observed isolate count sits on the boundary of its
support, so method of moments can match it only as
\(\theta \to -\infty\); the same network sampled to permit entry and
exit puts the statistic in the interior and the effect estimates
cleanly.\label{fig:boundary}}
\end{figure}

The second is a silent-failure pattern rather than a theorem, and it bit
this package. \texttt{RSiena::includeEffects} matches effect names by
non-standard evaluation, so passing a name held in a variable silently
selects nothing. A version of \textbf{searchnet} passed
\texttt{shortName = eff\$effect} into \texttt{...}, where the
unevaluated expression was deparsed and no effect matched; because the
call was wrapped in
\texttt{tryCatch(..., error = function(e) warning(...))}, the simulation
continued with the covariate effect \emph{absent}. Every actor-covariate
effect was affected, so results labelled ``heterogeneous strategies''
were in fact homogeneous-actor results. The package now ships a
regression test asserting that every effect declared in a
\texttt{structure\_model} appears with \texttt{include == TRUE} in the
fitted effects object. The design lesson generalises beyond this
package: \texttt{tryCatch}-to-warning is right for interactive
exploration and wrong for a simulation engine whose output gets
published.

\section{Summary and future directions}\label{sec:conclusion}

\textbf{searchnet} provides an \textsf{R} package that bridges two
established research traditions---NK fitness landscapes from strategy
and organization theory, and stochastic actor-oriented models from
social network analysis---through a bipartite network reformulation of
the search environment. The package makes three contributions to the
computational social science toolkit.

First, it \textbf{endogenizes the fitness landscape}. By representing
the search environment as a bipartite network governed by a SAOM,
landscape ruggedness becomes an emergent property of actors' affiliation
decisions rather than an exogenously imposed parameter. Researchers
studying adaptation, competitive dynamics, or organizational design can
now model settings where the landscape itself is shaped by the actors
who navigate it.

Second, it \textbf{decomposes \(K\)} into a four-dimensional coupled
degree system. The \(\{K\}\) framework reveals that the conventional
complexity parameter is one dimension of a structurally coupled process
linking actor behavior to landscape topology. The framework does not
introduce new network statistics; its contribution is interpretive,
providing a theoretical bridge between the NK and bipartite network
literatures that makes visible the transmission mechanism through which
strategic choices propagate into environmental structure.

Third, it provides a \textbf{complete simulation and analysis toolkit}.
The functional API (\texttt{saomnk\_env}, \texttt{saomnk\_model},
\texttt{saomnk\_run}) lowers the barrier to entry for management
scholars unfamiliar with \textbf{RSiena} syntax. The \(\{K\}\)-4 panel,
3D phase space visualization, utility decomposition, network snapshots,
and specification robustness diagnostics equip researchers with the
tools needed to design, execute, and validate computational experiments.
The package ships with eight vignettes, 16 Manim animation scenes, a
100-step formal proof chain (Parts A--L, including Theorem 5's
Brock--Durlauf recovery), 300 unit tests, and interactive game and
classroom modules that make the framework accessible to students and
practitioners.

Several directions for future development are planned. First,
integration with \textbf{RSiena}'s estimation machinery will allow
researchers to fit the SaoMNK model to observed bipartite network panel
data, enabling empirical calibration of simulation parameters. Second,
the fitness landscape module will be extended with scalable
approximation methods (e.g., Walsh decomposition, landscape correlation
analysis) to support problem sizes beyond \(N = 20\). Third, a
\textbf{Shiny} dashboard is under development to provide an interactive
visual interface for simulation design and exploration, lowering the
adoption barrier for researchers without programming experience. Fourth,
a web API via \textbf{plumber} will enable integration with external
frontends and cross-platform deployment. Two initially planned future
directions have already been realized in the current release: the
strategy game mode (Section\textasciitilde{}\ref{sec:interactive})
provides gamified engagement with the framework, and the classroom
teaching module offers a Capsim-style instructional platform with
industry presets, AI opponents, and automated debrief generation. The
accompanying PDW workshop vignette
(\texttt{vignette("saomnk-pdw-workshop")}) provides a complete
three-hour curriculum for introducing the framework to strategy and
organization scholars.

A companion paper \citep{downing2026architecture} applies the SaoMNK
framework empirically, using eight data layers from U.S.~airline
competitive dynamics to demonstrate how the \(\{K\}\) decomposition
reveals structural mechanisms invisible to the classical scalar \(K\)
parameterization.

\textbf{searchnet} is available from GitHub at
\url{https://github.com/sdownin/searchnet} and is under preparation for
CRAN submission. The package requires \textsf{R} \(\geq\) 4.1.0 and
depends on \textbf{R6}, \textbf{RSiena}, \textbf{igraph},
\textbf{ggplot2}, \textbf{dplyr}, \textbf{Matrix}, \textbf{network}, and
\textbf{reshape2}.

\section{Computational details}\label{sec:compdetails}

The results in this paper were obtained using
\textsf{R}\textasciitilde4.3.2 with the packages
\textbf{searchnet}\textasciitilde0.1.0,
\textbf{RSiena}\textasciitilde1.4.7,
\textbf{igraph}\textasciitilde2.0.3,
\textbf{ggplot2}\textasciitilde3.5.0, and
\textbf{Matrix}\textasciitilde1.6-5. \textsf{R} itself and all packages
used are available from the Comprehensive \textsf{R} Archive Network
(CRAN) at \url{https://CRAN.R-project.org/}.

\begin{Shaded}
\begin{Highlighting}[]
\FunctionTok{sessionInfo}\NormalTok{()}
\CommentTok{\#\textgreater{} R version 4.5.3 (2026{-}03{-}11 ucrt)}
\CommentTok{\#\textgreater{} Platform: x86\_64{-}w64{-}mingw32/x64}
\CommentTok{\#\textgreater{} Running under: Windows 11 x64 (build 26200)}
\CommentTok{\#\textgreater{} }
\CommentTok{\#\textgreater{} Matrix products: default}
\CommentTok{\#\textgreater{}   LAPACK version 3.12.1}
\CommentTok{\#\textgreater{} }
\CommentTok{\#\textgreater{} locale:}
\CommentTok{\#\textgreater{} [1] LC\_COLLATE=English\_United States.utf8 }
\CommentTok{\#\textgreater{} [2] LC\_CTYPE=English\_United States.utf8   }
\CommentTok{\#\textgreater{} [3] LC\_MONETARY=English\_United States.utf8}
\CommentTok{\#\textgreater{} [4] LC\_NUMERIC=C                          }
\CommentTok{\#\textgreater{} [5] LC\_TIME=English\_United States.utf8    }
\CommentTok{\#\textgreater{} }
\CommentTok{\#\textgreater{} time zone: America/Chicago}
\CommentTok{\#\textgreater{} tzcode source: internal}
\CommentTok{\#\textgreater{} }
\CommentTok{\#\textgreater{} attached base packages:}
\CommentTok{\#\textgreater{} [1] stats     graphics  grDevices utils     datasets  methods   base     }
\CommentTok{\#\textgreater{} }
\CommentTok{\#\textgreater{} other attached packages:}
\CommentTok{\#\textgreater{} [1] Matrix\_1.7{-}4    ggplot2\_4.0.2   searchnet\_0.9.2 rmarkdown\_2.31 }
\CommentTok{\#\textgreater{} }
\CommentTok{\#\textgreater{} loaded via a namespace (and not attached):}
\CommentTok{\#\textgreater{}   [1] mnormt\_2.1.2          pROC\_1.19.0.1         DRDID\_1.3.0          }
\CommentTok{\#\textgreater{}   [4] gridExtra\_2.3         rlang\_1.1.7           magrittr\_2.0.4       }
\CommentTok{\#\textgreater{}   [7] dreamerr\_1.5.0        otel\_0.2.0            compiler\_4.5.3       }
\CommentTok{\#\textgreater{}  [10] mgcv\_1.9{-}4            vctrs\_0.7.1           reshape2\_1.4.5       }
\CommentTok{\#\textgreater{}  [13] stringr\_1.6.0         pkgconfig\_2.0.3       fastmap\_1.2.0        }
\CommentTok{\#\textgreater{}  [16] backports\_1.5.0       labeling\_0.4.3        ggraph\_2.2.2         }
\CommentTok{\#\textgreater{}  [19] prodlim\_2025.04.28    network\_1.20.0        purrr\_1.2.1          }
\CommentTok{\#\textgreater{}  [22] xfun\_0.56             cachem\_1.1.0          trust\_0.1{-}9          }
\CommentTok{\#\textgreater{}  [25] jsonlite\_2.0.0        stringmagic\_1.2.0     recipes\_1.3.1        }
\CommentTok{\#\textgreater{}  [28] fastglm\_0.0.4         uuid\_1.2{-}2            tweenr\_2.0.3         }
\CommentTok{\#\textgreater{}  [31] psych\_2.6.5           broom\_1.0.12          parallel\_4.5.3       }
\CommentTok{\#\textgreater{}  [34] R6\_2.6.1              stringi\_1.8.7         RColorBrewer\_1.1{-}3   }
\CommentTok{\#\textgreater{}  [37] parallelly\_1.46.1     car\_3.1{-}5             rpart\_4.1.24         }
\CommentTok{\#\textgreater{}  [40] lubridate\_1.9.5       Rcpp\_1.1.1            BMisc\_1.4.8          }
\CommentTok{\#\textgreater{}  [43] iterators\_1.0.14      knitr\_1.51            future.apply\_1.20.2  }
\CommentTok{\#\textgreater{}  [46] RSiena\_1.5.0          timechange\_0.4.0      splines\_4.5.3        }
\CommentTok{\#\textgreater{}  [49] nnet\_7.3{-}20           igraph\_2.2.2          tidyselect\_1.2.1     }
\CommentTok{\#\textgreater{}  [52] abind\_1.4{-}8           yaml\_2.3.12           viridis\_0.6.5        }
\CommentTok{\#\textgreater{}  [55] timeDate\_4052.112     codetools\_0.2{-}20      listenv\_0.10.1       }
\CommentTok{\#\textgreater{}  [58] lattice\_0.22{-}9        tibble\_3.3.1          plyr\_1.8.9           }
\CommentTok{\#\textgreater{}  [61] withr\_3.0.2           S7\_0.2.1              coda\_0.19{-}4.1        }
\CommentTok{\#\textgreater{}  [64] evaluate\_1.0.5        future\_1.69.0         survival\_3.8{-}6       }
\CommentTok{\#\textgreater{}  [67] polyclip\_1.10{-}7       texreg\_1.39.5         pillar\_1.11.1        }
\CommentTok{\#\textgreater{}  [70] ggpubr\_0.6.3          carData\_3.0{-}6         stats4\_4.5.3         }
\CommentTok{\#\textgreater{}  [73] foreach\_1.5.2         generics\_0.1.4        scales\_1.4.0         }
\CommentTok{\#\textgreater{}  [76] globals\_0.19.0        class\_7.3{-}23          glue\_1.8.0           }
\CommentTok{\#\textgreater{}  [79] did\_2.5.0             tools\_4.5.3           data.table\_1.18.2.1  }
\CommentTok{\#\textgreater{}  [82] ModelMetrics\_1.2.2.2  gower\_1.0.2           ggsignif\_0.6.4       }
\CommentTok{\#\textgreater{}  [85] forcats\_1.0.1         visNetwork\_2.1.4      graphlayouts\_1.2.3   }
\CommentTok{\#\textgreater{}  [88] tidygraph\_1.3.1       cowplot\_1.2.0         grid\_4.5.3           }
\CommentTok{\#\textgreater{}  [91] bigmemory\_4.6.4       tidyr\_1.3.2           ipred\_0.9{-}15         }
\CommentTok{\#\textgreater{}  [94] nlme\_3.1{-}168          ggforce\_0.5.0         Formula\_1.2{-}5        }
\CommentTok{\#\textgreater{}  [97] cli\_3.6.5             bigmemory.sri\_0.1.8   viridisLite\_0.4.3    }
\CommentTok{\#\textgreater{} [100] lava\_1.8.2            dplyr\_1.2.0           gtable\_0.3.6         }
\CommentTok{\#\textgreater{} [103] rstatix\_0.7.3         digest\_0.6.39         caret\_7.0{-}1          }
\CommentTok{\#\textgreater{} [106] ggrepel\_0.9.7         htmlwidgets\_1.6.4     farver\_2.1.2         }
\CommentTok{\#\textgreater{} [109] memoise\_2.0.1         htmltools\_0.5.9       lifecycle\_1.0.5      }
\CommentTok{\#\textgreater{} [112] hardhat\_1.4.2         httr\_1.4.8            statnet.common\_4.13.0}
\CommentTok{\#\textgreater{} [115] MASS\_7.3{-}65}
\end{Highlighting}
\end{Shaded}


\renewcommand\refname{References}
\bibliography{searchnet.bib}

\end{document}
